\documentclass[11pt]{article}

\usepackage{amsmath,amssymb,amsthm}
\usepackage{hyperref}
\usepackage{xurl}
\usepackage{microtype}
\usepackage{xcolor}
\usepackage{enumitem}
\usepackage{geometry}
\usepackage{graphicx}
\usepackage{cite}
\usepackage{multirow}
\usepackage{array}
\geometry{margin=1in}
\emergencystretch=3em

% Theorem styles
\newtheoremstyle{definitionstyle}
  {3pt}{3pt}{\normalfont}{}{\bfseries}{.}{.5em}{}
\theoremstyle{definitionstyle}
\newtheorem{definition}{Definition}[section]
\newtheorem{theorem}{Theorem}[section]
\newtheorem{lemma}[theorem]{Lemma}
\newtheorem{corollary}[theorem]{Corollary}
\newtheorem{proposition}[theorem]{Proposition}
\newtheorem{assumption}{Assumption}[section]

\theoremstyle{plain}
\newtheorem*{remark}{Remark}

% Game environment
\newenvironment{game}[2][]
  {\par\addvspace{\topsep}\noindent\textbf{Game #2}\ #1\par\nobreak\smallskip\hrule\smallskip}
  {\hrule\par\addvspace{\topsep}}

\newcommand{\advantage}[1]{\mathbf{Adv}^{\mathrm{#1}}}
\newcommand{\negl}{\mathrm{negl}}
\newcommand{\bit}{\{0,1\}}
\newcommand{\DB}{\mathcal{DB}}

\title{Epoch-Bound Usage Tokens: Rate-Limited Anonymous Access with Zero-Knowledge Enforcement}
\author{Anant Singh}
\date{}

\begin{document}

\maketitle

\begin{abstract}
We present \emph{Epoch-Bound Usage Tokens (EBUT)}, a privacy-preserving primitive for enforcing rate-limited access in anonymous systems. In settings where users interact without persistent identities, enforcing usage policies such as rate limits and bounded resource consumption is challenging, as traditional approaches rely on identity-based tracking.

EBUT addresses this problem by introducing a token-based mechanism that combines epoch-bound issuance with bounded usage enforcement. Each user derives at most one token per epoch via a deterministic nullifier construction, preventing token accumulation across time and ensuring strict rate limits. Within an epoch, tokens support incremental usage through a recursive spend-and-refresh mechanism, enabling fine-grained control over resource consumption.

The protocol is realized using blind BBS+ signatures, Pedersen commitments, and zero-knowledge proofs, allowing the server to verify compliance with usage policies without learning any information about the user's identity, token linkage, or activity patterns. We formalize the security properties of EBUT and prove that it achieves unforgeability, unlinkability, and rate-limit soundness under standard cryptographic assumptions, including the $q$-Strong Diffie--Hellman, Discrete Logarithm, and SXDH assumptions in the random oracle model. To achieve practical efficiency, we incorporate a cross-curve batched equality proof that combines BLS12-381 and Ristretto groups, enabling fast Bulletproofs while maintaining rigorous security. We provide a complete, production-ready specification with detailed implementation guidance, API schemas, and performance estimates.
\end{abstract}

\section{Introduction}

Privacy-preserving systems aim to enable users to access services without revealing their identity or linking their actions across time. While modern cryptographic constructions provide strong anonymity guarantees, they struggle to enforce usage policies such as rate limits and bounded resource consumption. This tension arises because traditional enforcement mechanisms rely on persistent identities, which inherently violate unlinkability.

This problem becomes particularly acute in systems where resources are costly or constrained. For example, in decentralized storage or subscription-based services, a provider must ensure that each user consumes only a limited amount of resources over time. However, in an anonymous setting, the server cannot distinguish between users or track their behavior across sessions. As a result, a malicious user may accumulate tokens or credentials over time and later consume resources in a burst, violating rate limits and potentially causing service disruption.

Existing anonymous token systems partially address this challenge. Mechanisms such as Anonymous Credit Tokens (ACT)~\cite{act} enable users to obtain tokens and spend them in a privacy-preserving manner, supporting features such as partial spending and unlinkable redemption. However, these systems fundamentally treat tokens as accumulable resources: users may obtain tokens over time and spend them arbitrarily later. Consequently, they do not enforce temporal constraints on usage and cannot prevent burst consumption.

A natural approach is to introduce explicit rate-limiting mechanisms, such as per-epoch quotas enforced via zero-knowledge proofs. However, naïve constructions based on accumulators or membership proofs require maintaining dynamic sets and performing complex updates, leading to significant inefficiencies and poor usability. Moreover, such approaches often introduce additional trust assumptions or degrade privacy guarantees.

In this work, we propose a different approach: instead of enforcing rate limits externally, we embed temporal constraints directly into the structure of tokens themselves. We introduce \emph{Epoch-Bound Usage Tokens (EBUT)}, a primitive in which tokens are intrinsically bound to discrete time epochs. Each user can derive at most one token per epoch through a deterministic nullifier construction, ensuring that token issuance is inherently rate-limited and preventing accumulation across time.

Within each epoch, EBUT supports fine-grained usage through a recursive spend-and-refresh mechanism, allowing tokens to be consumed incrementally while preserving unlinkability. All correctness conditions are enforced via zero-knowledge proofs, enabling the server to verify compliance without learning any information about the user's identity or activity patterns.

The key insight of our construction is that temporal structure can be used as a cryptographic enforcement mechanism. By binding tokens to epochs at the protocol level, EBUT eliminates the need for external rate-limiting logic while preserving strong anonymity guarantees.

\subsection{Our Contributions}
We summarize our main contributions as follows:
\begin{itemize}[leftmargin=*]
\item \textbf{A new primitive for rate-limited anonymous access.} We introduce \emph{Epoch-Bound Usage Tokens (EBUT)}, a cryptographic primitive that enforces rate limits in anonymous systems. Unlike existing anonymous token mechanisms, EBUT prevents accumulation of tokens across time by binding issuance to discrete epochs.
\item \textbf{Temporal enforcement via deterministic nullifiers.} We propose a deterministic epoch-based nullifier construction of the form $N_T = k \cdot H_T$, which ensures that each user can derive at most one token per epoch while preserving unlinkability across epochs. This embeds rate-limiting constraints directly into the cryptographic structure of tokens, eliminating the need for external enforcement mechanisms.
\item \textbf{Non-accumulating token model.} We formalize and realize a token model in which resources cannot be accumulated over time. This directly addresses the burst-consumption problem present in existing anonymous token systems, where users may hoard tokens and later spend them arbitrarily.
\item \textbf{Recursive intra-epoch usage mechanism.} We design a recursive spend-and-refresh protocol that enables fine-grained resource consumption within a single epoch. This allows tokens to be used incrementally while maintaining strict global rate limits.
\item \textbf{Formalization of rate-limit soundness.} We introduce and formalize the notion of \emph{rate-limit soundness} for anonymous token systems, capturing the requirement that no adversary can exceed the allowed rate of token generation or usage over time.
\item \textbf{Efficient zero-knowledge enforcement.} We construct EBUT using blind BBS+ signatures, Pedersen commitments, and zero-knowledge proofs, enabling the server to verify policy compliance without learning any information about the user's identity or activity. The protocol achieves practical efficiency with proof sizes on the order of a few kilobytes and minimal server-side state.
\item \textbf{Security proofs under standard assumptions.} We prove that EBUT achieves unforgeability, unlinkability, and rate-limit soundness under standard cryptographic assumptions, including the $q$-Strong Diffie--Hellman assumption, the Discrete Logarithm assumption, and the SXDH assumption in the random oracle model.
\item \textbf{Cross-curve batched equality proofs.} We provide a complete, self-contained treatment of a dual-curve range proof system that combines BLS12-381 and Ristretto groups, with full security proofs (completeness, knowledge soundness, zero-knowledge). This is of independent interest for protocols requiring efficient range proofs on committed values across different curves.
\end{itemize}

\subsection{Related Work}
\paragraph{Anonymous credentials.} Anonymous credential systems enable users to authenticate without revealing their identity or linking multiple interactions. Classical constructions, including those based on pairing-based signatures such as BBS+, provide unlinkable credential presentation and selective disclosure~\cite{CamenischL01,CL02,BBS04}. These systems focus primarily on identity privacy and attribute verification, but do not directly address enforcement of usage policies such as rate limits or bounded resource consumption.

\paragraph{Anonymous tokens and usage-based systems.} Several works extend anonymous credentials to support token-based access and usage control. In particular, Anonymous Credit Tokens (ACT)~\cite{act} introduce a model in which users obtain tokens that can be spent incrementally in a privacy-preserving manner. Such systems support partial spending and unlinkable redemption, making them suitable for usage-based billing and resource allocation. However, these constructions fundamentally treat tokens as accumulable resources; a user may obtain tokens over time and later spend them arbitrarily, which allows burst consumption and violates temporal rate limits.

\paragraph{Rate limiting in anonymous systems.} Rate limiting in the absence of identities is inherently challenging. Existing approaches often rely on external mechanisms such as per-epoch quotas enforced via accumulators, Merkle trees, or membership proofs combined with zero-knowledge protocols~\cite{BonehShoup}. While these approaches can enforce constraints, they typically require maintaining dynamic state, performing frequent updates, or introducing additional trust assumptions. Moreover, they may incur significant computational overhead and degrade usability. Other approaches use pseudonymous identifiers or linkable tags to enforce limits, but these weaken privacy guarantees by introducing linkability across sessions or epochs.

\paragraph{Temporal enforcement mechanisms.} Some systems incorporate time-based validity, such as expiring credentials or time-limited tokens. However, these mechanisms restrict \emph{when} a token can be used, rather than \emph{how frequently} tokens can be generated or consumed. They do not prevent accumulation across time and therefore do not enforce strict rate limits.

\paragraph{Non-transferable anonymous tokens.} The recent work of Durak et al.~\cite{NTAT} introduces the notion of non-transferable anonymous tokens, where a token is bound to a client's secret key and cannot be transferred without revealing that key. While their focus is on preventing token sharing, our work addresses the orthogonal problem of rate limiting across time. Interestingly, both approaches use a deterministic nullifier to enforce a policy, but in EBUT the nullifier is epoch-based rather than session-based.

\paragraph{Our approach.} In contrast to prior work, EBUT embeds temporal constraints directly into the token structure. By binding token issuance to discrete epochs via a deterministic nullifier construction, EBUT ensures that each user can derive at most one token per epoch. This eliminates token accumulation across time while preserving unlinkability. Furthermore, EBUT combines this temporal enforcement with a recursive spend-and-refresh mechanism, enabling fine-grained usage within each epoch without sacrificing privacy. This provides a unified solution to rate limiting in anonymous systems, without requiring external stateful enforcement mechanisms or weakening anonymity guarantees.

\section{Preliminaries}
We denote by $\lambda \in \mathbb{N}$ the security parameter. For a finite set $S$, we write $x \xleftarrow{\$} S$ to indicate that $x$ is chosen uniformly at random from $S$. A function $\varepsilon(\lambda)$ is \emph{negligible} if $\varepsilon(\lambda) = \lambda^{-\omega(1)}$. We use $\mathsf{negl}(\lambda)$ to denote an arbitrary negligible function.

\subsection{Bilinear Groups and Cryptographic Assumptions}
Let $\mathbb{G}_1, \mathbb{G}_2, \mathbb{G}_T$ be cyclic groups of prime order $q$ (depending on $\lambda$) with generators $g_1 \in \mathbb{G}_1$, $g_2 \in \mathbb{G}_2$. A bilinear map (pairing) is an efficiently computable map $e: \mathbb{G}_1 \times \mathbb{G}_2 \to \mathbb{G}_T$ satisfying:
\begin{itemize}
\item Bilinearity: $e(g_1^a, g_2^b) = e(g_1, g_2)^{ab}$ for all $a,b \in \mathbb{Z}_q$.
\item Non-degeneracy: $e(g_1, g_2) \neq 1_{\mathbb{G}_T}$.
\end{itemize}
We assume a Type-3 pairing where no efficiently computable homomorphism exists between $\mathbb{G}_1$ and $\mathbb{G}_2$.

\begin{assumption}[Discrete Logarithm (DL)]
Given $(g, g^x)$ where $g$ is a generator of a cyclic group $\mathbb{G}$ of prime order $q$ and $x \xleftarrow{\$} \mathbb{Z}_q$, no probabilistic polynomial-time (PPT) adversary can output $x$ with more than negligible probability.
\end{assumption}

\begin{assumption}[$q$-Strong Diffie--Hellman ($q$-SDH)~\cite{BB04}]
Let $\mathbb{G}_1, \mathbb{G}_2$ be groups of prime order $q$ with generators $g_1, g_2$. Given a $(q+2)$-tuple $(g_1, g_2, g_2^\gamma, g_2^{\gamma^2}, \dots, g_2^{\gamma^q})$ for random $\gamma \xleftarrow{\$} \mathbb{Z}_q^*$, no PPT adversary can output a pair $(x, g_1^{1/(\gamma+x)})$ with $x \in \mathbb{Z}_q \setminus \{-\gamma\}$, except with negligible probability.
\end{assumption}

\begin{assumption}[Symmetric eXternal Diffie--Hellman (SXDH)]
In a Type-3 pairing setting, the Decisional Diffie--Hellman (DDH) assumption holds in both $\mathbb{G}_1$ and $\mathbb{G}_2$. That is, given $(g_1, g_1^a, g_1^b, g_1^c)$, it is hard to decide whether $c = ab$ or random; similarly in $\mathbb{G}_2$.
\end{assumption}

\subsection{BBS+ Signatures}
We use the BBS+ signature scheme~\cite{BBS04,CFRG-BBS} over a bilinear group setting. The scheme allows signing a vector of messages $(m_1,\dots,m_L) \in \mathbb{Z}_q^L$.

\begin{itemize}
\item $\mathsf{Setup}(1^\lambda)$: Output parameters $(\mathbb{G}_1,\mathbb{G}_2,\mathbb{G}_T,e,q,g_1,g_2)$ and additional generators $h_0,\dots,h_L \in \mathbb{G}_1$, derived via a hash-to-curve function with domain separation.
\item $\mathsf{KeyGen}()$: Choose $sk \xleftarrow{\$} \mathbb{Z}_q$, set $pk = g_2^{sk}$.
\item $\mathsf{Sign}(sk, (m_1,\dots,m_L))$: Choose random $e, s \xleftarrow{\$} \mathbb{Z}_q$ and compute
  \[
  A = \left( g_1 \cdot h_0^s \cdot \prod_{i=1}^L h_i^{m_i} \right)^{\frac{1}{e+sk}} \in \mathbb{G}_1.
  \]
  Output signature $(A,e,s)$.
\item $\mathsf{Verify}(pk, (m_1,\dots,m_L), (A,e,s))$: Check $A \neq 1_{\mathbb{G}_1}$, $e \neq 0$, and
  \[
  e(A, pk \cdot g_2^e) = e\left( g_1 \cdot h_0^s \cdot \prod_{i=1}^L h_i^{m_i}, g_2 \right).
  \]
\end{itemize}
BBS+ signatures are existentially unforgeable under chosen-message attacks (EUF-CMA) under the $q$-SDH assumption.

\subsection{Zero-Knowledge Proofs}
We employ non-interactive zero-knowledge proofs of knowledge (NIZKPoK) for various relations, constructed via the Fiat--Shamir transform in the random oracle model. Specifically, we use:
\begin{itemize}
\item A proof of possession of a BBS+ signature (Protocol~1 in \cite{act}), which allows proving knowledge of a valid signature on a set of hidden messages, with optional selective disclosure.
\item Schnorr-style proofs for equality of discrete logarithms across different bases.
\item Bulletproofs~\cite{Bulletproofs} for range proofs, adapted to be bound to the outer proof context.
\end{itemize}
All proofs include a domain separator and a context binding to the application to prevent cross-protocol attacks.

\subsection{Time Epochs and Deterministic Nullifiers}
Time is divided into discrete epochs $T \in \{0,1,\dots\}$ of fixed duration. We define a public deterministic hash-to-curve function
\[
H_T = \mathcal{H}_{\mathbb{G}_1}(\text{``EBUT:Epoch:''} \parallel T) \in \mathbb{G}_1.
\]
For a user secret $k \in \mathbb{Z}_q^*$, the epoch nullifier is $N_T = k \cdot H_T$. Under the SXDH assumption, the values $N_T$ for different epochs are computationally unlinkable.

\section{Syntax and Security Definitions}
\label{sec:defs}

We now formalize the syntax and security requirements for an Epoch-Bound Usage Token (EBUT) scheme. The protocol involves a server $\mathcal{S}$ and multiple clients $\mathcal{C}$. The server is responsible for issuing tokens and verifying their usage; clients request tokens and spend them to access resources. The scheme operates over discrete time epochs $T = 0,1,2,\dots$, and each client is associated with a long-term master secret key.

\subsection{Syntax}
An EBUT scheme consists of the following polynomial-time algorithms and interactive protocols.

\begin{itemize}[leftmargin=*]
\item \textbf{$\mathsf{Setup}(1^\lambda) \to \mathsf{pp}$.} On input the security parameter $\lambda$, outputs public parameters $\mathsf{pp}$ which are common to all parties.
\item \textbf{$\mathsf{Server.KeyGen}(\mathsf{pp}) \to (\mathsf{sk}_S, \mathsf{pk}_S)$.} Generates the server's long-term key pair.
\item \textbf{$\mathsf{Client.KeyGen}(\mathsf{pp}) \to (\mathsf{sk}_C, \mathsf{pk}_C)$.} Generates a client's master key pair. The secret key $\mathsf{sk}_C$ is a scalar $k_{\mathsf{sub}} \in \mathbb{Z}_q^*$, and the public key $\mathsf{pk}_C$ is a group element $g_1^{\mathsf{sk}_C}$. We assume a public-key infrastructure where the server knows the client's public key (or a commitment to it) during issuance.
\item \textbf{$\mathsf{Issue}$ (interactive).} An interactive protocol between a client and the server to obtain a master subscription token.
  \begin{itemize}
  \item $\mathsf{Client.Query}(\mathsf{pp}, \mathsf{sk}_C, \mathsf{pk}_S) \to (\mathsf{query}, \mathsf{st})$: The client produces a request for a master token and local state.
  \item $\mathsf{Server.Issue}(\mathsf{pp}, \mathsf{sk}_S, \mathsf{pk}_C, \mathsf{query}) \to \mathsf{resp}$: The server verifies the client's request and, if valid, returns a blind signature on the client's secret key and subscription parameters.
  \item $\mathsf{Client.Final}(\mathsf{st}, \mathsf{resp}) \to (\mathcal{T}_{\mathsf{master}}, \omega)$: The client completes the issuance, obtaining a master token $\mathcal{T}_{\mathsf{master}}$ and a decommitment $\omega$.
  \end{itemize}
  The master token $\mathcal{T}_{\mathsf{master}}$ contains a BBS+ signature on $(k_{\mathsf{sub}}, c_{\mathsf{max}}, E_{\mathsf{max}})$, where $c_{\mathsf{max}}$ is the maximum credit per epoch and $E_{\mathsf{max}}$ is the subscription expiry epoch.
\item \textbf{$\mathsf{Refresh}$ (interactive).} An interactive protocol executed at the beginning of an epoch $T \le E_{\mathsf{max}}$ to derive an epoch-bound daily token.
  \begin{itemize}
  \item $\mathsf{Client.Refresh}(\mathsf{pp}, \mathcal{T}_{\mathsf{master}}, T) \to (\mathsf{req}_{\mathsf{ref}}, \mathsf{st})$: The client produces a refresh request including a zero-knowledge proof of a valid master token, the epoch nullifier $N_T$, and a commitment to a fresh daily secret $k_{\mathsf{daily}}$.
  \item $\mathsf{Server.VerifyRefresh}(\mathsf{pp}, \mathsf{sk}_S, T, \mathsf{req}_{\mathsf{ref}}) \to \mathsf{resp}_{\mathsf{ref}} / \bot$: The server verifies the proof, checks that the epoch nullifier $N_T$ is fresh (not in $\DB_{\mathsf{epoch}}$), and issues a blind BBS+ signature on $(c_{\mathsf{max}}, k_{\mathsf{daily}}, T)$.
  \item $\mathsf{Client.FinalRefresh}(\mathsf{st}, \mathsf{resp}_{\mathsf{ref}}) \to \mathcal{T}_{\mathsf{daily}}$: The client finalizes the daily token $\mathcal{T}_{\mathsf{daily}}$ for epoch $T$.
  \end{itemize}
\item \textbf{$\mathsf{Spend}$ (interactive).} An interactive protocol, parameterised by a mode $\mu \in \{\mathsf{Exact},\mathsf{Hidden}\}$, to spend $s$ credits from a daily token and obtain a refund token for the remaining balance $m = c_{\mathsf{bal}} - s$.
  \begin{itemize}
  \item $\mathsf{Client.Spend}(\mathsf{pp}, \mathcal{T}_{\mathsf{current}}, \mu, \mathsf{aux}) \to (\mathsf{req}_{\mathsf{spend}}, \mathsf{st})$: In Exact mode, the auxiliary input is a required amount $s_{\mathsf{req}}$ supplied by the application and the client proves the disclosed spend equals it. In Hidden mode, the client chooses $s \in [1,c_{\mathsf{bal}}]$ locally and proves the relation in zero-knowledge.
  \item $\mathsf{Server.VerifySpend}(\mathsf{pp}, \mathsf{sk}_S, \mu, \mathsf{aux}, \mathsf{req}_{\mathsf{spend}}) \to (\mathsf{resp}_{\mathsf{spend}}, \mathsf{out})$: The server verifies the mode-specific proof, checks that the token's nullifier $k_{\mathsf{cur}}$ is fresh (not in $\DB_{\mathsf{spend}}$), and issues a blind refund token. In Exact mode the server additionally enforces $s = s_{\mathsf{req}}$. Output $\mathsf{out} = 1$ if accepted, $0$ otherwise.
  \item $\mathsf{Client.FinalSpend}(\mathsf{st}, \mathsf{resp}_{\mathsf{spend}}) \to \mathcal{T}_{\mathsf{new}}$: The client finalizes the new token $\mathcal{T}_{\mathsf{new}}$ with balance $m$.
  \end{itemize}
\end{itemize}

The server maintains two databases for double-spending prevention:
\begin{itemize}
\item $\DB_{\mathsf{epoch}}$: stores epoch nullifiers $N_T \in \mathbb{G}_1$ observed during $\mathsf{Refresh}$; entries carry a TTL equal to the epoch duration plus the grace period, and are expired individually rather than bulk-cleared at the epoch boundary.
\item $\DB_{\mathsf{spend}}$: stores token nullifiers $k_{\mathsf{cur}} \in \mathbb{Z}_q^*$ observed during $\mathsf{Spend}$; entries carry a TTL equal to the epoch duration plus the grace period, and are expired individually rather than bulk-cleared at the epoch boundary.
\end{itemize}
Additionally, the server may employ an idempotency cache to handle retries gracefully.

\subsection{Correctness}
An EBUT scheme must satisfy \emph{correctness}: for any honest execution of the protocols, the server accepts and the client obtains valid tokens with the expected balances. Formally, for all $\lambda$, all $\mathsf{pp} \leftarrow \mathsf{Setup}(1^\lambda)$, all key pairs $(\mathsf{sk}_S, \mathsf{pk}_S)$, $(\mathsf{sk}_C, \mathsf{pk}_C)$, and all valid parameters $(c_{\mathsf{max}}, E_{\mathsf{max}})$:
\begin{itemize}
\item If $\mathsf{Issue}$ runs honestly, the client outputs a master token $\mathcal{T}_{\mathsf{master}}$ that verifies under $\mathsf{pk}_S$.
\item For any epoch $T \le E_{\mathsf{max}}$, an honest $\mathsf{Refresh}$ yields a daily token $\mathcal{T}_{\mathsf{daily}}$ with balance $c_{\mathsf{max}}$ that is accepted by the server.
\item For any spend amount $1 \le s \le c_{\mathsf{bal}}$ and either spend mode $\mu \in \{\mathsf{Exact},\mathsf{Hidden}\}$ (with $s = s_{\mathsf{req}}$ in Exact mode), an honest $\mathsf{Spend}$ yields a new token with balance $c_{\mathsf{bal}}-s$ and the server outputs $\mathsf{out}=1$.
\end{itemize}
Correctness follows directly from the completeness of the underlying zero-knowledge proofs and the BBS+ signature scheme.

\subsection{Security Definitions}
We define three core security properties: \emph{unforgeability}, \emph{unlinkability}, and \emph{rate-limit soundness}. The first two are standard for anonymous token systems; the third captures the epoch-bound rate limiting specific to EBUT.

\subsubsection{Unforgeability}
Unforgeability ensures that an adversary cannot obtain more valid tokens than it is entitled to. We formulate this as a one-more unforgeability game where the adversary can interact with the server via issuance, refresh, and spend oracles. The goal is to produce more tokens (or spend more credits) than the number of master tokens issued.

\begin{game}{1 (Unforgeability, $\mathsf{UNF}$)}
Let $\mathcal{A}$ be a PPT adversary. The game proceeds as follows:
\begin{enumerate}[leftmargin=*]
\item $\mathsf{pp} \leftarrow \mathsf{Setup}(1^\lambda)$, $(\mathsf{sk}_S, \mathsf{pk}_S) \leftarrow \mathsf{Server.KeyGen}(\mathsf{pp})$.
\item $\mathcal{A}$ is given $\mathsf{pp}$ and $\mathsf{pk}_S$, and has access to the following oracles:
  \begin{itemize}
  \item $\mathcal{O}_{\mathsf{Issue}}(\mathsf{pk}_C)$: Runs the $\mathsf{Issue}$ protocol with $\mathcal{A}$ acting as the client. The server uses $\mathsf{pk}_C$ as the client's public key. Increments a counter $q_{\mathsf{issue}}[\mathsf{pk}_C]$.
  \item $\mathcal{O}_{\mathsf{Refresh}}(T, \mathsf{pk}_C, \pi)$: $\mathcal{A}$ supplies a refresh request for epoch $T$ and public key $\mathsf{pk}_C$. The server verifies and, if valid, returns a daily token. Increments $q_{\mathsf{refresh}}[\mathsf{pk}_C][T]$.
  \item $\mathcal{O}_{\mathsf{Spend}}(s, \mathsf{req})$: $\mathcal{A}$ supplies a spend request. The server verifies and, if valid, returns a refund token. Increments $q_{\mathsf{spend}}$.
  \end{itemize}
\item $\mathcal{A}$ outputs a tuple $(\mathsf{pk}_C^*, \mathcal{T}_1, \dots, \mathcal{T}_k, \omega_1, \dots, \omega_k)$ and an epoch $T^*$.
\item $\mathcal{A}$ wins if for the target public key $\mathsf{pk}_C^*$:
  \begin{itemize}
  \item All tokens $\mathcal{T}_i$ are valid for epoch $T^*$ (or later) and their total credit value exceeds $c_{\mathsf{max}} \cdot q_{\mathsf{issue}}[\mathsf{pk}_C^*]$; \emph{or}
  \item The number of valid daily tokens obtained for $\mathsf{pk}_C^*$ across all epochs exceeds $q_{\mathsf{issue}}[\mathsf{pk}_C^*]$; \emph{or}
  \item There exists an epoch $T \le E_{\mathsf{max}}$ where the number of valid daily tokens for $\mathsf{pk}_C^*$ exceeds $1$ (i.e., more than one refresh per epoch).
  \end{itemize}
\end{enumerate}
The advantage of $\mathcal{A}$ is $\advantage{UNF}(\mathcal{A}) = \Pr[\mathcal{A} \text{ wins}]$. An EBUT scheme is \emph{unforgeable} if for all PPT $\mathcal{A}$, $\advantage{UNF}(\mathcal{A}) \le \negl(\lambda)$.
\end{game}

\noindent \textbf{Remark.} The oracle queries capture both malicious clients and honest-but-curious server interactions. The game directly models the rate-limiting property: even with adaptive issuance queries, the adversary cannot exceed the per-epoch token limit or the total credit bound.

\subsubsection{Unlinkability}
Unlinkability requires that the server cannot link a refresh or spend transaction to a specific master token or to other transactions of the same client, beyond what is revealed by the public nullifiers (which are unlinkable across epochs).

We define a game where the adversary (acting as a malicious server) chooses two clients and must distinguish which one is participating in a challenge session.

\begin{game}{2 (Unlinkability, $\mathsf{UNLINK}$)}
Let $\mathcal{A} = (\mathcal{A}_1, \mathcal{A}_2)$ be a PPT adversary. The game proceeds as follows:
\begin{enumerate}[leftmargin=*]
\item $\mathsf{pp} \leftarrow \mathsf{Setup}(1^\lambda)$, $(\mathsf{sk}_S, \mathsf{pk}_S) \leftarrow \mathsf{Server.KeyGen}(\mathsf{pp})$. $\mathcal{A}$ is given $\mathsf{pp}$ and $\mathsf{pk}_S$.
\item $\mathcal{A}_1$ outputs two client key pairs $(\mathsf{sk}_C^0, \mathsf{pk}_C^0)$ and $(\mathsf{sk}_C^1, \mathsf{pk}_C^1)$, along with state $\mathsf{st}$. Both must be valid.
\item The challenger flips a fair coin $b \xleftarrow{\$} \{0,1\}$. 
\item $\mathcal{A}_2(\mathsf{st})$ interacts with the challenger in two phases:
  \begin{itemize}
  \item \textbf{Issuance phase:} For $i \in \{0,1\}$, the challenger runs $\mathsf{Issue}$ with $\mathcal{A}_2$ acting as the server, using client $i$'s secret key. The adversary learns the master tokens (if any) but not the internal randomness.
  \item \textbf{Challenge phase:} $\mathcal{A}_2$ selects an epoch $T$ and initiates a $\mathsf{Refresh}$ protocol. The challenger runs the client side using $\mathsf{sk}_C^b$ (the challenge client). The adversary observes the refresh request. Similarly, $\mathcal{A}_2$ may request a $\mathsf{Spend}$ protocol for a token owned by $b$.
  \end{itemize}
\item $\mathcal{A}_2$ outputs a guess $b' \in \{0,1\}$.
\end{enumerate}
$\mathcal{A}$ wins if $b' = b$. The advantage is $\advantage{UNLINK}(\mathcal{A}) = |\Pr[b' = b] - 1/2|$. An EBUT scheme is \emph{unlinkable} if for all PPT $\mathcal{A}$, $\advantage{UNLINK}(\mathcal{A}) \le \negl(\lambda)$.
\end{game}

\noindent The definition captures that even a malicious server controlling the issuance process cannot link the refresh/spend of a client to its master token or to other epochs.

\subsubsection{Rate-Limit Soundness}
This new notion formalizes the epoch-bound rate limiting: no adversary can obtain more than one daily token per epoch per master secret, and cannot spend more than $c_{\mathsf{max}}$ credits within an epoch. We define a game where the adversary controls both the server and clients, but must respect the token generation limits.

\begin{game}{3 (Rate-Limit Soundness, $\mathsf{RLS}$)}
Let $\mathcal{A}$ be a PPT adversary. The game proceeds as follows:
\begin{enumerate}[leftmargin=*]
\item $\mathsf{pp} \leftarrow \mathsf{Setup}(1^\lambda)$. $\mathcal{A}$ may generate server keys (possibly maliciously).
\item $\mathcal{A}$ outputs a tuple $(\mathsf{pk}_S, T, \{(N_T^{(i)}, \pi_{\mathsf{refresh}}^{(i)}, \mathcal{T}_{\mathsf{daily}}^{(i)})\}_{i=1}^n)$ where each $\mathcal{T}_{\mathsf{daily}}^{(i)}$ is a daily token for epoch $T$ under $\mathsf{pk}_S$, and $\pi_{\mathsf{refresh}}^{(i)}$ is the corresponding refresh proof.
\item $\mathcal{A}$ wins if:
  \begin{itemize}
  \item All daily tokens are valid (i.e., $\mathsf{Server.VerifyRefresh}$ accepts each).
  \item The epoch nullifiers $N_T^{(i)}$ are all distinct and non-zero.
  \item There exist $i \neq j$ such that the underlying master secret $k_{\mathsf{sub}}$ is the same (i.e., the tokens were derived from the same master token), yet the adversary obtained $n \ge 2$ daily tokens for the same epoch.
  \end{itemize}
  \emph{or} the adversary produces a set of spend transactions within epoch $T$ that consume more than $c_{\mathsf{max}}$ total credits from a single master token without double-spending the same token nullifier.
\end{enumerate}
The advantage $\advantage{RLS}(\mathcal{A}) = \Pr[\mathcal{A} \text{ wins}]$. An EBUT scheme satisfies \emph{rate-limit soundness} if for all PPT $\mathcal{A}$, $\advantage{RLS}(\mathcal{A}) \le \negl(\lambda)$.
\end{game}

\noindent \textbf{Discussion.} Rate-limit soundness directly captures the core guarantee of EBUT: it is infeasible to bypass the per-epoch token limit. This is stronger than simple unforgeability because even if the adversary obtains a valid master token, it cannot use it to exceed the epoch quota. The security reduction will rely on the binding properties of the zero-knowledge proofs and the collision-resistance of the nullifier construction.

\section{Construction of Epoch-Bound Usage Tokens}
\label{sec:construction}

In this section we present a concrete instantiation of EBUT using blind BBS+ signatures, Pedersen commitments, and zero-knowledge proofs. The construction relies on a bilinear group setting (Type-3 pairing) and the random oracle model.

\subsection{Public Parameters and Key Generation}
\label{sec:keygen}

Let $\mathbb{G}_1, \mathbb{G}_2, \mathbb{G}_T$ be groups of prime order $q$ with a pairing $e: \mathbb{G}_1 \times \mathbb{G}_2 \to \mathbb{G}_T$. Let $g_1 \in \mathbb{G}_1$ and $g_2 \in \mathbb{G}_2$ be fixed generators.

\paragraph{Public Parameters.}
The global parameters $\mathsf{pp}$ consist of:
\begin{itemize}
\item The bilinear group description $(\mathbb{G}_1,\mathbb{G}_2,\mathbb{G}_T,e,q,g_1,g_2)$.
\item Six independent generators $h_0, h_1, h_2, h_3, h_4, h_5 \in \mathbb{G}_1$ for BBS+ signatures and Pedersen commitments. Here $h_4$ is used for balance commitments and $h_5$ is reserved for hidden spend-amount commitments. These are derived via a hash-to-curve function $\mathcal{H}_{\mathbb{G}_1}$ with domain separation, e.g., $h_i = \mathcal{H}_{\mathbb{G}_1}(\text{``EBUT:Generator:''} \parallel i)$.
\item A context hash $H_{\mathsf{ctx}}$ that binds the protocol instance to an application identifier, the server public keys, and all public generators, preventing cross-deployment replay and parameter-substitution attacks. Specifically, the production formula is:
  \[
  \begin{aligned}
  H_{\mathsf{ctx}} = \mathcal{H}_{\mathsf{scalar}}\!\bigl(
    &\text{``ACT:Context''} \parallel \mathsf{AppID}
    \parallel \mathsf{pk}_{\mathsf{master}} \parallel \mathsf{pk}_{\mathsf{daily}} \\
    &\parallel g_1 \parallel g_2
    \parallel h_0 \parallel h_1 \parallel h_2 \parallel h_3 \parallel h_4 \parallel h_5
  \bigr).
  \end{aligned}
  \]
  where $\mathcal{H}_{\mathsf{scalar}}: \{0,1\}^* \to \mathbb{Z}_q$ is a hash function modeled as a random oracle. Binding the generators prevents a parameter-substitution attack in which an adversary convinces a client to use generators for which the adversary knows discrete-logarithm relations; binding both server public keys (master and daily, see Section~\ref{sec:implementation-triage}) prevents cross-key replay. All elements are serialised in compressed form before hashing. \textbf{Implementation note on hash width:} For the BLS12-381 scalar field ($q$ a 255-bit prime), reducing a 32-byte SHA-256 digest modulo $q$ introduces a per-value bias of $(2^{256} \bmod q)/q \approx 2^{-255}$, which is negligible at the 128-bit security level. The reference implementation uses SHA-256 with repeated-subtraction mod-$q$ reduction and is cryptographically sound. Implementations targeting conservative interoperability may use SHA-512 (64-byte output) before reduction as a belt-and-suspenders measure.
\end{itemize}

\paragraph{Server Key Generation.}
The server generates \emph{two} independent key pairs: a master signing key $(\mathsf{sk}_{\mathsf{master}}, \mathsf{pk}_{\mathsf{master}}) = (y_m, g_2^{y_m})$ used to sign master subscription tokens, and a daily signing key $(\mathsf{sk}_{\mathsf{daily}}, \mathsf{pk}_{\mathsf{daily}}) = (y_d, g_2^{y_d})$ used to sign daily tokens. Using a single key pair is sufficient for the security proofs but inadvisable in practice: key separation ensures that a compromise of the daily signing key (e.g., via a hot-path HSM breach) does not invalidate long-term master subscriptions, and the daily key can be rotated at epoch boundaries. See Section~\ref{sec:implementation-triage} for further discussion.

To enable the Unlinkability proof (Theorem~\ref{thm:unlinkability}) against a malicious server, the $\mathsf{Server.KeyGen}$ algorithm is extended to produce a Non-Interactive Zero-Knowledge Proof of Knowledge (NIZKPoK) of the secret key. For each key $y \in \{y_m, y_d\}$ and corresponding $\mathsf{pk}_S = g_2^y$, the server produces a proof $\pi_S = (T_y, z_y)$ as follows:
\begin{enumerate}[leftmargin=*,nosep]
  \item Choose $\rho_y \xleftarrow{\$} \mathbb{Z}_q$ and compute $T_y = g_2^{\rho_y} \in \mathbb{G}_2$.
  \item Compute the Fiat--Shamir challenge:
    \[
    c_y = \mathcal{H}_{\mathsf{scalar}}\!\bigl(\text{``EBUT:ServerPoK''} \parallel g_2 \parallel \mathsf{pk}_S \parallel T_y\bigr).
    \]
  \item Compute the response $z_y = \rho_y + c_y \cdot y \pmod{q}$.
  \item Output $\pi_S = (T_y, z_y)$.
\end{enumerate}
Clients and simulators \textbf{must} explicitly reject any $\mathsf{pk}_S$ for which $g_2^{z_y} \neq T_y \cdot \mathsf{pk}_S^{c_y}$. This proof injects a cryptographic trapdoor into the Setup phase: the standard rewindable extractor on $\mathcal{H}_{\mathsf{scalar}}$ allows the security-proof Simulator to extract $y$ from $\pi_S$ and simulate BBS+ transcripts without knowledge of the client's master secret (see Game~$\mathsf{H}_2$ in Section~\ref{sec:unlinkability}).

\paragraph{Client Key Generation.}
Each client generates a master secret key $k_{\mathsf{sub}} \xleftarrow{\$} \mathbb{Z}_q^*$ (non-zero) and computes the corresponding public key $X = g_1^{k_{\mathsf{sub}}} \in \mathbb{G}_1$. The server is assumed to know $X$ (or a commitment to it) out-of-band during the initial subscription setup.

\paragraph{Grace Period.}
The server MAY accept refresh or spend requests for epoch $T-1$ during a configurable grace window $\Delta_{\mathsf{grace}} \ge 0$ after the start of epoch $T$. If $\Delta_{\mathsf{grace}} = 0$ (the default), requests for expired epochs are immediately rejected. When $\Delta_{\mathsf{grace}} > 0$, the nullifier databases must retain entries for at least $\Delta_{\mathsf{grace}}$ beyond the epoch boundary (see Section~\ref{sec:implementation}).

\subsection{Master Subscription Token Issuance}
The issuance protocol allows a client to obtain a master token that encodes the subscription limits: maximum credits per epoch $c_{\mathsf{max}} \in [0, 2^{32}-1]$ and expiry epoch $E_{\mathsf{max}}$. The protocol is blind: the server does not learn the client's secret $k_{\mathsf{sub}}$, only a commitment to it.

\begin{enumerate}[leftmargin=*]
\item \textbf{Client commitment.}
  The client chooses randomness $r_{\mathsf{sub}} \xleftarrow{\$} \mathbb{Z}_q$ and computes a Pedersen commitment to $k_{\mathsf{sub}}$:
  \[
  K_{\mathsf{sub}} = k_{\mathsf{sub}} \cdot h_1 + r_{\mathsf{sub}} \cdot h_0 \in \mathbb{G}_1.
  \]
  The client sends $K_{\mathsf{sub}}$ to the server along with a Schnorr proof of knowledge $\pi_{\mathsf{PoK}}$ of the opening $(k_{\mathsf{sub}}, r_{\mathsf{sub}})$.

  The proof $\pi_{\mathsf{PoK}}$ is a standard Schnorr protocol: the prover picks $\alpha, \beta \xleftarrow{\$} \mathbb{Z}_q$, computes $T = \alpha \cdot h_1 + \beta \cdot h_0$, obtains challenge $c = \mathcal{H}_{\mathsf{scalar}}(H_{\mathsf{ctx}} \parallel K_{\mathsf{sub}} \parallel T \parallel \text{"PoK"})$, and responds with $s_k = \alpha + c \cdot k_{\mathsf{sub}}$, $s_r = \beta + c \cdot r_{\mathsf{sub}}$. The verifier checks $s_k \cdot h_1 + s_r \cdot h_0 \stackrel{?}{=} T + c \cdot K_{\mathsf{sub}}$.

\item \textbf{Server signing.}
  The server knows the subscription parameters $c_{\mathsf{max}}$ and $E_{\mathsf{max}}$ (which are policy-dependent and can be public). It chooses fresh randomness $e_{\mathsf{sub}}, s'_{\mathsf{sub}} \xleftarrow{\$} \mathbb{Z}_q$ and computes a BBS+ signature on the message vector $(k_{\mathsf{sub}}, c_{\mathsf{max}}, E_{\mathsf{max}})$ with bases $(h_1, h_2, h_3)$:
  \[
  A_{\mathsf{sub}} = \left( g_1 \cdot K_{\mathsf{sub}} \cdot h_0^{s'_{\mathsf{sub}}} \cdot h_2^{c_{\mathsf{max}}} \cdot h_3^{E_{\mathsf{max}}} \right)^{\frac{1}{e_{\mathsf{sub}} + y}} \in \mathbb{G}_1.
  \]
  The server sends $(A_{\mathsf{sub}}, e_{\mathsf{sub}}, s'_{\mathsf{sub}})$ to the client.

\item \textbf{Client finalisation.}
  The client sets $s_{\mathsf{sub}} = r_{\mathsf{sub}} + s'_{\mathsf{sub}} \bmod q$. The master token is
  \[
  \mathcal{T}_{\mathsf{master}} = (A_{\mathsf{sub}}, e_{\mathsf{sub}}, s_{\mathsf{sub}}, k_{\mathsf{sub}}, c_{\mathsf{max}}, E_{\mathsf{max}}).
  \]
  The client verifies the signature using the BBS+ verification equation with $\mathsf{pk}_S$.
\end{enumerate}

\subsection{Epoch Refresh: Obtaining a Daily Token}
At the beginning of epoch $T \le E_{\mathsf{max}}$, the client derives an epoch-bound daily token. The process uses a zero-knowledge proof that the client holds a valid master token and that the epoch nullifier $N_T$ is correctly computed.

\subsubsection{Client-Side Computation}
\begin{enumerate}[leftmargin=*]
\item \textbf{Epoch nullifier.}
  Compute the epoch base point $H_T = \mathcal{H}_{\mathbb{G}_1}(\text{``EBUT:Epoch:''} \parallel T) \in \mathbb{G}_1$.
  The epoch nullifier is $N_T = k_{\mathsf{sub}} \cdot H_T \in \mathbb{G}_1$.
\item \textbf{Daily token commitment.}
  Derive the daily secret and randomness deterministically from the master secret and epoch $T$ via a PRF modelled as a Random Oracle:
  \[
  k_{\mathsf{daily}} = \mathcal{H}_{\mathsf{scalar}}\!\bigl(\text{``EBUT:Derive:K''} \parallel k_{\mathsf{sub}} \parallel T\bigr), \quad
  r_{\mathsf{daily}} = \mathcal{H}_{\mathsf{scalar}}\!\bigl(\text{``EBUT:Derive:R''} \parallel k_{\mathsf{sub}} \parallel T\bigr).
  \]
  Form the commitment:
  \[
  K_{\mathsf{daily}} = c_{\mathsf{max}} \cdot h_4 + k_{\mathsf{daily}} \cdot h_1 + T \cdot h_2 + r_{\mathsf{daily}} \cdot h_0 \in \mathbb{G}_1.
  \]
  This commitment binds the daily token to epoch $T$, its credit balance $c_{\mathsf{max}}$, and the fresh secret $k_{\mathsf{daily}}$.

  \textbf{Security note.}
  Because $k_{\mathsf{sub}}$ has min-entropy $\ge 250$ bits (DL is hard given $X = g_1^{k_{\mathsf{sub}}}$), and $\mathcal{H}_{\mathsf{scalar}}$ is modelled as a Random Oracle, the outputs $k_{\mathsf{daily}}$ and $r_{\mathsf{daily}}$ are computationally indistinguishable from uniform in $\mathbb{Z}_q$ to any PPT adversary.  The Pedersen commitment $K_{\mathsf{daily}}$ therefore retains perfect hiding.

  \textbf{Liveness guarantee.}
  If the client crashes after the server burns $N_T$ in $\DB_{\mathsf{epoch}}$ but before the response arrives, the client need only re-evaluate $N_T = k_{\mathsf{sub}} \cdot H_T$ from the persisted $k_{\mathsf{sub}}$ and the public epoch $T$.  The server's idempotency cache (keyed on $N_T$; see Section~\ref{sec:idempotency}) returns the cached blind signature, which the client unblinds correctly.  No synchronous disk write of ephemeral randomness is required; the client may re-derive $k_{\mathsf{daily}}$ and $r_{\mathsf{daily}}$ fresh.

\item \textbf{Expiry range proof (Cross-Curve BEQ).}
  Compute $\Delta = E_{\mathsf{max}} - T$.
  Choose $r_{\Delta} \xleftarrow{\$} \mathbb{Z}_{q_{\text{BLS}}}$ and $r_{25519}^{(\Delta)} \xleftarrow{\$} \mathbb{Z}_{q_{25519}}$.
  Compute BLS12-381 commitment: $C_{\Delta} = \Delta \cdot h_3 + r_{\Delta} \cdot h_0$.
  Compute Ristretto commitment: $C_{25519}^{(\Delta)} = \Delta \cdot G + r_{25519}^{(\Delta)} \cdot H$.
  Generate the BEQ proof $\pi_{\text{BEQ\_exp}}$ for statement $(C_{\Delta}, C_{25519}^{(\Delta)})$ with auxiliary commitments $(A', \bar{A}, T_{\mathsf{BBS}}, N_T, K_{\mathsf{daily}})$ (the BBS+ proof components are defined below).
  Define a bridge commitment $C_{\mathsf{bridge}} = C_{\Delta} + T \cdot h_3 = E_{\mathsf{max}} \cdot h_3 + r_{\Delta} \cdot h_0$.
\end{enumerate}

\subsubsection[Zero-Knowledge Proof pi refresh]{Zero-Knowledge Proof $\pi_{\mathsf{refresh}}$}
The client proves knowledge of a valid master token and correct derivation of $N_T$, $K_{\mathsf{daily}}$, and $C_{\mathsf{bridge}}$. The proof is a variant of the BBS+ zero-knowledge proof of possession (Protocol~1 in~\cite{act}) with additional Schnorr equalities.

\begin{enumerate}[leftmargin=*]
\item \textbf{BBS+ randomization.}
  Pick $r_1 \xleftarrow{\$} \mathbb{Z}_q^*$ and compute
  \[
  A' = A_{\mathsf{sub}}^{r_1}, \quad
  \bar{A} = A'^{-e_{\mathsf{sub}}} \cdot \left( g_1 \cdot h_0^{s_{\mathsf{sub}}} \cdot h_1^{k_{\mathsf{sub}}} \cdot h_2^{c_{\mathsf{max}}} \cdot h_3^{E_{\mathsf{max}}} \right)^{r_1}.
  \]
  Define scaled secrets: $\tilde{k} = k_{\mathsf{sub}} r_1$, $\tilde{c} = c_{\mathsf{max}} r_1$, $\tilde{E} = E_{\mathsf{max}} r_1$, $\tilde{s} = s_{\mathsf{sub}} r_1$.
\item \textbf{Commitments.}
  Choose random blinders $\rho_e, \rho_{r_1}, \rho_{\tilde{s}}, \rho_{\tilde{k}}, \rho_{\tilde{c}}, \rho_{\tilde{E}} \xleftarrow{\$} \mathbb{Z}_q$ and compute
  \[
  T_{\mathsf{BBS}} = A'^{-\rho_e} \cdot g_1^{\rho_{r_1}} \cdot h_0^{\rho_{\tilde{s}}} \cdot h_1^{\rho_{\tilde{k}}} \cdot h_2^{\rho_{\tilde{c}}} \cdot h_3^{\rho_{\tilde{E}}} \in \mathbb{G}_1.
  \]
  Additionally, choose blinders $\rho_u, \rho_v, \rho_w \xleftarrow{\$} \mathbb{Z}_q$ (for $k_{\mathsf{daily}} r_1$, $r_{\mathsf{daily}} r_1$, $r_{\Delta} r_1$).
  Compute bridging commitments:
  \begin{align*}
    T_{\mathsf{scale},N} &= \rho_{r_1} \cdot N_T, & T_N &= \rho_{\tilde{k}} \cdot H_T, \\
    T_{\mathsf{scale},D} &= \rho_{r_1} \cdot K_{\mathsf{daily}}, & T_D &= \rho_{\tilde{c}} \cdot h_4 + \rho_u \cdot h_1 + \rho_{r_1} \cdot T \cdot h_2 + \rho_v \cdot h_0, \\
    T_{\mathsf{scale},B} &= \rho_{r_1} \cdot C_{\mathsf{bridge}}, & T_B &= \rho_{\tilde{E}} \cdot h_3 + \rho_w \cdot h_0.
  \end{align*}
\item \textbf{Challenge.}
  Compute the Fiat--Shamir challenge as
  \[
  c = \mathcal{H}_{\mathsf{scalar}}\left(
  \begin{array}{c}
  H_{\mathsf{ctx}} \parallel \mathsf{pk}_S \parallel T \parallel N_T \parallel K_{\mathsf{daily}} \parallel C_{\Delta} \parallel \pi_{\text{BEQ\_exp}} \parallel \\
  A' \parallel \bar{A} \parallel T_{\mathsf{BBS}} \parallel T_{\mathsf{scale},N} \parallel T_N \parallel T_{\mathsf{scale},D} \parallel T_D \parallel T_{\mathsf{scale},B} \parallel T_B \parallel \text{``Refresh''}
  \end{array}
  \right).
  \]
\item \textbf{Responses.}
  Compute modulo $q$:
  \begin{align*}
    z_e &= \rho_e + c \cdot e_{\mathsf{sub}}, & z_{r_1} &= \rho_{r_1} + c \cdot r_1, \\
    z_{\tilde{s}} &= \rho_{\tilde{s}} + c \cdot \tilde{s}, & z_{\tilde{k}} &= \rho_{\tilde{k}} + c \cdot \tilde{k}, \\
    z_{\tilde{c}} &= \rho_{\tilde{c}} + c \cdot \tilde{c}, & z_{\tilde{E}} &= \rho_{\tilde{E}} + c \cdot \tilde{E}, \\
    z_u &= \rho_u + c \cdot (k_{\mathsf{daily}} r_1), & z_v &= \rho_v + c \cdot (r_{\mathsf{daily}} r_1), \\
    z_w &= \rho_w + c \cdot (r_{\Delta} r_1).
  \end{align*}
\item \textbf{Output proof.}
  $\pi_{\mathsf{refresh}}$ consists of:
  \begin{itemize}[leftmargin=*, nosep]
  \item The randomized BBS+ components $A', \bar{A}, T_{\mathsf{BBS}}$.
  \item The nullifier bridging commitments $T_{\mathsf{scale},N}, T_N$.
  \item The daily token bridging commitments $T_{\mathsf{scale},D}, T_D$.
  \item The expiry bridging commitments $T_{\mathsf{scale},B}, T_B$.
  \item The scalar responses $z_e, z_{r_1}, z_{\tilde{s}}, z_{\tilde{k}}, z_{\tilde{c}}, z_{\tilde{E}}, z_u, z_v, z_w$.
  \item The cross-curve range proof $\pi_{\text{BEQ\_exp}}$.
  \end{itemize}
\end{enumerate}

\subsubsection{Server Verification}
The server performs the following checks in order:
\begin{enumerate}[leftmargin=*]
\item \textbf{Idempotency cache check (FIRST):} Compute
  \[
  \mathtt{key}_{\mathsf{ref}} = \mathcal{H}_{\mathsf{scalar}}\!\bigl(\text{EBUT:Idempotency:Refresh} \parallel N_T\bigr).
  \]
  The key is derived exclusively from $N_T$ (the epoch nullifier), which is uniquely determined by the client's master secret $k_{\mathsf{sub}}$ and the public epoch $T$ via $N_T = k_{\mathsf{sub}} \cdot H_T$.  An honest client reconstructing the same $N_T$ deterministically will always hit the same cache slot.  \textbf{Critically, the key must not be derived from $\pi_{\mathsf{refresh}}$ or $K_{\mathsf{daily}}$ alone}: those values depend on the optional blinding randomness $r_{\mathsf{daily}}$, which an adversarial client may vary across retries while presenting the same epoch nullifier, enabling a double-mint if a cache miss follows a partial write (see Section~\ref{sec:idempotency}).  If $\mathtt{key}_{\mathsf{ref}}$ is present in the idempotency cache, return the cached $(A_{\mathsf{daily}}, e_{\mathsf{daily}}, s'_{\mathsf{daily}})$ immediately \textbf{without} re-executing the expensive verification or touching $\DB_{\mathsf{epoch}}$.
\item \textbf{Epoch check:} $T$ must equal the current epoch (or be within a small grace window).
\item \textbf{Nullifier freshness check (read-only):} $N_T \notin \DB_{\mathsf{epoch}}$ and $N_T \neq 1_{\mathbb{G}_1}$. \textbf{Do not insert $N_T$ at this step.} Inserting before all proofs are verified creates a TOCTOU vulnerability: an adversary can submit a valid $N_T$ with garbage ZK proofs; the server would prematurely burn the epoch nullifier and then abort during proof verification, permanently locking the legitimate user out for the entire epoch. All state mutations are deferred to the atomic step~11.
\item \textbf{Basic bounds:} $A' \neq 1_{\mathbb{G}_1}$ and $T_{\mathsf{BBS}} \neq 1_{\mathbb{G}_1}$.
\item \textbf{Expiry bridge:} $C_{\mathsf{bridge}} \stackrel{?}{=} C_{\Delta} + T \cdot h_3$.
\item \textbf{Recompute challenge $c$} from the received values (including $H_{\mathsf{ctx}}$ and $\mathsf{pk}_S$).
\item \textbf{Verify bridging equalities (subtraction form):}
  \begin{align*}
    z_{\tilde{k}} \cdot H_T - T_N &\stackrel{?}{=} z_{r_1} \cdot N_T - T_{\mathsf{scale},N}, \\
    z_{\tilde{c}} \cdot h_4 + z_u \cdot h_1 + z_{r_1} \cdot T \cdot h_2 + z_v \cdot h_0 - T_D &\stackrel{?}{=} z_{r_1} \cdot K_{\mathsf{daily}} - T_{\mathsf{scale},D}, \\
    z_{\tilde{E}} \cdot h_3 + z_w \cdot h_0 - T_B &\stackrel{?}{=} z_{r_1} \cdot C_{\mathsf{bridge}} - T_{\mathsf{scale},B}.
  \end{align*}
\item \textbf{Schnorr multi-scalar multiplication (BBS+ core):}
  \[
  A'^{-z_e} \cdot g_1^{z_{r_1}} \cdot h_0^{z_{\tilde{s}}} \cdot h_1^{z_{\tilde{k}}} \cdot h_2^{z_{\tilde{c}}} \cdot h_3^{z_{\tilde{E}}} \stackrel{?}{=} T_{\mathsf{BBS}} \cdot \bar{A}^c.
  \]
\item \textbf{BEQ verification:} Verify $\pi_{\text{BEQ\_exp}}$ with respect to $C_{\Delta}$.
\item \textbf{Pairing check:}
  \[
  e(A', \mathsf{pk}_S) \stackrel{?}{=} e(\bar{A}, g_2).
  \]
\item \textbf{Atomic state mutation (step 11, LAST):} Only after all checks in steps~1--10 pass: first write $\mathtt{key}_{\mathsf{ref}} \mapsto (A_{\mathsf{daily}}, e_{\mathsf{daily}}, s'_{\mathsf{daily}})$ to the idempotency cache (where $\mathtt{key}_{\mathsf{ref}} = \mathcal{H}_{\mathsf{scalar}}(\text{EBUT:Idempotency:Refresh} \parallel N_T)$), then atomically insert $N_T$ into $\DB_{\mathsf{epoch}}$ (reject if already present, indicating a concurrent duplicate or replay). \textbf{The idempotency cache entry must be persisted before the $\DB_{\mathsf{epoch}}$ insertion.} Because the cache key is derived from $N_T$, any retry that presents the same epoch nullifier---regardless of whether $K_{\mathsf{daily}}$ or $r_{\mathsf{daily}}$ was regenerated---will hit the existing cache entry and receive the original response, rather than bypassing it and triggering a second mint.  If both writes are to separate storage systems, they must be wrapped in a single distributed transaction or the cache must be written first; writing to $\DB_{\mathsf{epoch}}$ first voids the idempotency guarantee and will permanently lock out a client whose network dropped between the two writes.
\end{enumerate}

If all checks pass, the server issues a daily token by choosing $e_{\mathsf{daily}}, s'_{\mathsf{daily}} \xleftarrow{\$} \mathbb{Z}_q$ and computing
\[
A_{\mathsf{daily}} = \left( g_1 \cdot K_{\mathsf{daily}} \cdot h_0^{s'_{\mathsf{daily}}} \right)^{\frac{1}{e_{\mathsf{daily}} + y}}.
\]
It sends $(A_{\mathsf{daily}}, e_{\mathsf{daily}}, s'_{\mathsf{daily}})$ to the client (the response was already cached in the idempotency store during the atomic mutation step above).

\subsubsection{Client Finalisation}
The client sets $s_{\mathsf{daily}} = r_{\mathsf{daily}} + s'_{\mathsf{daily}} \bmod q$. The daily token is
\[
\mathcal{T}_{\mathsf{daily}} = (A_{\mathsf{daily}}, e_{\mathsf{daily}}, s_{\mathsf{daily}}, k_{\mathsf{daily}}, c_{\mathsf{max}}, T).
\]

\subsection{Intra-Epoch Spending}

Within epoch $T_{\mathsf{issue}}$, a client may spend $s$ credits from a current token
\[
\mathcal{T}_{\mathsf{current}} = (A_{\mathsf{cur}}, e_{\mathsf{cur}}, s_{\mathsf{cur}}, k_{\mathsf{cur}}, c_{\mathsf{bal}}, T_{\mathsf{issue}})
\]
and obtain a refund token for the remaining balance $m = c_{\mathsf{bal}} - s$.
The protocol supports two spend modes.

\paragraph{Mode Selection and Security Rationale.}
\emph{Exact mode} is used when the server or application endpoint enforces a fixed price $s_{\mathsf{req}}$ (for example, ``one API call costs exactly 3 units''). It is the simpler, smaller, and benchmarked wire format. \emph{Hidden mode} is used when the client chooses the spend amount locally and that amount should remain hidden from the server; it adds an additional commitment $C_s$, an additional BEQ proof, and extra bridge terms, and is therefore strictly larger and more expensive than Exact mode. A request \textbf{must commit to exactly one mode}: the mode tag and all mode-specific public fields are bound into both the Fiat--Shamir challenge and the idempotency key.

\subsubsection{Server-Enforced Exact Amount}
\label{sec:spend-exact}

The server supplies an application-level required amount $s_{\mathsf{req}} \in [1, c_{\mathsf{bal}}]$. The client sets $s = s_{\mathsf{req}}$ and discloses both values in the request. The server \textbf{must} reject unless $s = s_{\mathsf{req}}$ and all Hidden-mode-only fields are absent.

\paragraph{Client-Side Computation.}
\begin{enumerate}[leftmargin=*]
\item \textbf{Session nonce.}
  The client generates a 128-bit session nonce locally using a CSPRNG:
  \[
  \eta \xleftarrow{\$} \{0,1\}^{128}.
  \]
  No network request to the server is needed.  The nonce $\eta$ is bound into the Fiat--Shamir challenge and serves as a replay-prevention token.  The server checks $\eta \notin \mathcal{N}_{\mathsf{spent}}$ as a read-only check before proof verification, then atomically marks $\eta$ as spent only after all cryptographic checks pass.  By the birthday paradox, the probability of an honest client generating a colliding $\eta$ within a single epoch is at most $q_S^2/2^{128}$, which is negligible for any polynomial number of spends $q_S$.
\item \textbf{Refund commitment.}
  Let $m = c_{\mathsf{bal}} - s$. Choose a fresh secret $k^* \xleftarrow{\$} \mathbb{Z}_q^*$ and randomness $r^* \xleftarrow{\$} \mathbb{Z}_q$. Compute the commitment for the new token:
  \[
  K' = m \cdot h_4 + k^* \cdot h_1 + T_{\mathsf{issue}} \cdot h_2 + r^* \cdot h_0 \in \mathbb{G}_1.
  \]
\item \textbf{Range proof (Cross-Curve BEQ).}
  Choose $r_{\mathsf{BP}} \xleftarrow{\$} \mathbb{Z}_{q_{\text{BLS}}}$ and $r_{25519}^{(m)} \xleftarrow{\$} \mathbb{Z}_{q_{25519}}$.
  Compute BLS12-381 commitment: $C_{\mathsf{BP}} = m \cdot h_4 + r_{\mathsf{BP}} \cdot h_0$.
  Compute Ristretto commitment: $C_{25519}^{(m)} = m \cdot G + r_{25519}^{(m)} \cdot H$.
  Generate the BEQ proof $\pi_{\text{BEQ\_spend}}$ for statement $(C_{\mathsf{BP}}, C_{25519}^{(m)})$ with auxiliary commitments $(A', \bar{A}, T_{\mathsf{BBS}}, K', C_{\mathsf{total}})$ (the BBS+ proof components are defined below).
\item \textbf{Spend bridge.}
  \[
  C_{\mathsf{total}} = K' + s \cdot h_4 = c_{\mathsf{bal}} \cdot h_4 + k^* \cdot h_1 + T_{\mathsf{issue}} \cdot h_2 + r^* \cdot h_0.
  \]
\end{enumerate}

\paragraph{Zero-Knowledge Proof $\pi_{\mathsf{spend}}$ (Exact Mode).}
The client proves knowledge of a valid daily token on $(c_{\mathsf{bal}}, k_{\mathsf{cur}}, T_{\mathsf{issue}})$ with $k_{\mathsf{cur}}$ and $T_{\mathsf{issue}}$ disclosed, and that $m = c_{\mathsf{bal}} - s$ is correctly formed.

\begin{enumerate}[leftmargin=*]
\item \textbf{BBS+ randomization.}
  Pick $r_1 \xleftarrow{\$} \mathbb{Z}_q^*$ and compute
  \[
  A' = A_{\mathsf{cur}}^{r_1}, \quad
  \bar{A} = A'^{-e_{\mathsf{cur}}} \cdot \left( g_1 \cdot h_0^{s_{\mathsf{cur}}} \cdot h_4^{c_{\mathsf{bal}}} \cdot h_1^{k_{\mathsf{cur}}} \cdot h_2^{T_{\mathsf{issue}}} \right)^{r_1}.
  \]
  Define scaled secrets $\tilde{c} = c_{\mathsf{bal}} r_1$ and $\tilde{s} = s_{\mathsf{cur}} r_1$.
\item \textbf{Commitments.}
  Choose $\rho_e, \rho_{r_1}, \rho_{\tilde{s}}, \rho_{\tilde{c}} \xleftarrow{\$} \mathbb{Z}_q$ and compute
  \[
  T_{\mathsf{BBS}} = A'^{-\rho_e} \cdot g_1^{\rho_{r_1}} \cdot h_0^{\rho_{\tilde{s}}} \cdot h_4^{\rho_{\tilde{c}}} \cdot h_1^{k_{\mathsf{cur}} \cdot \rho_{r_1}} \cdot h_2^{T_{\mathsf{issue}} \cdot \rho_{r_1}}.
  \]
  Choose additional blinders $\rho_u, \rho_v, \rho_w \xleftarrow{\$} \mathbb{Z}_q$ (for $k^* r_1$, $r^* r_1$, $r_{\mathsf{BP}} r_1$). Define $\rho_m = \rho_{\tilde{c}} - s \cdot \rho_{r_1}$ and compute the bridge commitments:
  \begin{align*}
    T_{\mathsf{scale},R} &= \rho_{r_1} \cdot K', & T_{\mathsf{refund}} &= \rho_m \cdot h_4 + \rho_u \cdot h_1 + \rho_{r_1} \cdot T_{\mathsf{issue}} \cdot h_2 + \rho_v \cdot h_0, \\
    T_{\mathsf{scale},T} &= \rho_{r_1} \cdot C_{\mathsf{total}}, & T_{\mathsf{total}} &= \rho_{\tilde{c}} \cdot h_4 + \rho_u \cdot h_1 + \rho_{r_1} \cdot T_{\mathsf{issue}} \cdot h_2 + \rho_v \cdot h_0, \\
    T_{\mathsf{scale},BP} &= \rho_{r_1} \cdot C_{\mathsf{BP}}, & T_{\mathsf{BP}} &= \rho_m \cdot h_4 + \rho_w \cdot h_0.
  \end{align*}
\item \textbf{Challenge.}
  The challenge includes the mode tag, the server-enforced amount, and the replay-prevention nonce $\eta$:
  \[
  \begin{aligned}
  c = \mathcal{H}_{\mathsf{scalar}}\!\bigl(
    &H_{\mathsf{ctx}} \parallel \mathsf{pk}_S \parallel \text{``Exact''}
    \parallel s \parallel s_{\mathsf{req}} \parallel k_{\mathsf{cur}} \parallel T_{\mathsf{issue}} \\
    &\parallel \eta \parallel K' \parallel C_{\mathsf{total}} \parallel C_{\mathsf{BP}}
    \parallel \pi_{\text{BEQ\_spend}} \\
    &\parallel A' \parallel \bar{A} \parallel T_{\mathsf{BBS}}
    \parallel T_{\mathsf{scale},R} \parallel T_{\mathsf{refund}} \\
    &\parallel T_{\mathsf{scale},T} \parallel T_{\mathsf{total}}
    \parallel T_{\mathsf{scale},BP} \parallel T_{\mathsf{BP}} \parallel \text{``Spend''}
  \bigr).
  \end{aligned}
  \]
\item \textbf{Responses.}
  Compute:
  \begin{align*}
    z_e &= \rho_e + c \cdot e_{\mathsf{cur}}, & z_{r_1} &= \rho_{r_1} + c \cdot r_1, \\
    z_{\tilde{s}} &= \rho_{\tilde{s}} + c \cdot \tilde{s}, & z_{\tilde{c}} &= \rho_{\tilde{c}} + c \cdot \tilde{c}, \\
    z_u &= \rho_u + c \cdot (k^* r_1), & z_v &= \rho_v + c \cdot (r^* r_1), \\
    z_w &= \rho_w + c \cdot (r_{\mathsf{BP}} r_1).
  \end{align*}
  Define the derived response $z_m = z_{\tilde{c}} - s \cdot z_{r_1}$.
\item \textbf{Output proof.}
  $\pi_{\mathsf{spend}}$ consists of:
  \begin{itemize}[leftmargin=*, nosep]
  \item The randomized BBS+ components $A', \bar{A}, T_{\mathsf{BBS}}$.
  \item The refund bridging commitments $T_{\mathsf{scale},R}, T_{\mathsf{refund}}$.
  \item The spend-bridge commitments $T_{\mathsf{scale},T}, T_{\mathsf{total}}$.
  \item The range-proof bridging commitments $T_{\mathsf{scale},BP}, T_{\mathsf{BP}}$.
  \item The scalar responses $z_e, z_{r_1}, z_{\tilde{s}}, z_{\tilde{c}}, z_u, z_v, z_w$.
  \item The cross-curve range proof $\pi_{\text{BEQ\_spend}}$.
  \end{itemize}
\end{enumerate}

\paragraph{Server Verification (Exact Mode).}
\begin{enumerate}[leftmargin=*]
\item \textbf{Mode and amount checks.} Require $\mathsf{mode} = \text{``Exact''}$, require Exact-mode fields $(s, s_{\mathsf{req}})$ to be present, require Hidden-mode fields $(C_s, \pi_{\text{BEQ\_s}})$ to be absent, and reject unless $1 \le s = s_{\mathsf{req}} \le 2^{32}-1$.
\item \textbf{Epoch check.} The disclosed $T_{\mathsf{issue}}$ must equal the current epoch unless a non-zero grace window $\Delta_{\mathsf{grace}}$ is configured.
\item \textbf{Idempotency cache check (FIRST).} Compute $h_{\mathsf{id}}$ as
  the canonical spend cache key from Section~7.4 with
  $\mathsf{mode}=\text{``Exact''}$, with $s_{\mathsf{req}}^{\mathsf{enc}}$ set
  to the 4-byte encoding of $s_{\mathsf{req}}$, and with
  $C_s^{\mathsf{enc}} = \pi_{\text{BEQ\_s}}^{\mathsf{enc}} = \epsilon$.
  If $h_{\mathsf{id}}$ is in the idempotency cache, return the cached refund token immediately \textbf{without checking the double-spend database}.
\item \textbf{Nonce read-only check.} Verify that $\eta$ is non-zero and is exactly 16 bytes, and that $\eta \notin \mathcal{N}_{\mathsf{spent}}$. \textbf{Do not mark $\eta$ as spent at this step.}
\item \textbf{Double-spend read-only check.} Check that $k_{\mathsf{cur}} \notin \DB_{\mathsf{spend}}$, reject if $k_{\mathsf{cur}} = 0$, and \textbf{do not} insert into $\DB_{\mathsf{spend}}$ at this step.
\item \textbf{Basic bounds.} $A' \neq 1_{\mathbb{G}_1}$ and $T_{\mathsf{BBS}} \neq 1_{\mathbb{G}_1}$.
\item \textbf{Spend bridge.} Check $C_{\mathsf{total}} \stackrel{?}{=} K' + s \cdot h_4$.
\item \textbf{Recompute $c$.} Recompute the Fiat--Shamir challenge from the received values, including the mode tag and $s_{\mathsf{req}}$.
\item \textbf{Compute $z_m = z_{\tilde{c}} - s \cdot z_{r_1}$.}
\item \textbf{Verify bridging equalities.}
  \begin{align*}
    z_m \cdot h_4 + z_u \cdot h_1 + z_{r_1} \cdot T_{\mathsf{issue}} \cdot h_2 + z_v \cdot h_0 - T_{\mathsf{refund}} &\stackrel{?}{=} z_{r_1} \cdot K' - T_{\mathsf{scale},R}, \\
    z_{\tilde{c}} \cdot h_4 + z_u \cdot h_1 + z_{r_1} \cdot T_{\mathsf{issue}} \cdot h_2 + z_v \cdot h_0 - T_{\mathsf{total}} &\stackrel{?}{=} z_{r_1} \cdot C_{\mathsf{total}} - T_{\mathsf{scale},T}, \\
    z_m \cdot h_4 + z_w \cdot h_0 - T_{\mathsf{BP}} &\stackrel{?}{=} z_{r_1} \cdot C_{\mathsf{BP}} - T_{\mathsf{scale},BP}.
  \end{align*}
\item \textbf{Schnorr MSM.}
  \[
  A'^{-z_e} \cdot g_1^{z_{r_1}} \cdot h_0^{z_{\tilde{s}}} \cdot h_4^{z_{\tilde{c}}} \cdot h_1^{k_{\mathsf{cur}} \cdot z_{r_1}} \cdot h_2^{T_{\mathsf{issue}} \cdot z_{r_1}} \stackrel{?}{=} T_{\mathsf{BBS}} \cdot \bar{A}^c.
  \]
\item \textbf{BEQ verification.} Verify $\pi_{\text{BEQ\_spend}}$ with respect to $C_{\mathsf{BP}}$.
\item \textbf{Pairing check.} Check $e(A', \mathsf{pk}_S) \stackrel{?}{=} e(\bar{A}, g_2)$.
\item \textbf{Atomic state mutation (LAST).} Only after all checks pass: first write $h_{\mathsf{id}} \mapsto (A^*, e^*, s'^*)$ to the idempotency cache (computing $(A^*, e^*, s'^*)$ by signing $K'$ as described below), then atomically insert $k_{\mathsf{cur}}$ into $\DB_{\mathsf{spend}}$ and atomically mark $\eta$ as spent. \textbf{The idempotency cache entry must be written before the $\DB_{\mathsf{spend}}$ insertion.}
\end{enumerate}

If all checks pass and the state mutation succeeds, the server issues a refund token by choosing $e^*, s'^* \xleftarrow{\$} \mathbb{Z}_q$ and computing
\[
A^* = \left( g_1 \cdot K' \cdot h_0^{s'^*} \right)^{\frac{1}{e^* + y}}.
\]
It sends $(A^*, e^*, s'^*)$ to the client (the response was already cached during the atomic state-mutation step).

\paragraph{Client Finalisation (Exact Mode).}
The client sets $s^* = r^* + s'^* \bmod q$ and obtains the new token
\[
\mathcal{T}_{\mathsf{new}} = (A^*, e^*, s^*, k^*, m, T_{\mathsf{issue}}).
\]

\subsubsection{Client-Chosen Amount (Hidden)}
\label{sec:spend-hidden}

In Hidden mode the client chooses $s \in [1,c_{\mathsf{bal}}]$ locally and the server learns neither $s$ nor the refund balance $m = c_{\mathsf{bal}} - s$. Hidden mode reuses the Exact-mode refund token format but extends the proof with a separate commitment to the spend amount:
\[
C_s = s \cdot h_5 + r_s \cdot h_0 \in \mathbb{G}_1,
\]
where $h_5$ is an independent generator reserved for hidden spend amounts.

\paragraph{Client-Side Computation.}
The client computes $m = c_{\mathsf{bal}} - s$, forms the refund commitment
\[
K' = m \cdot h_4 + k^* \cdot h_1 + T_{\mathsf{issue}} \cdot h_2 + r^* \cdot h_0,
\]
and forms the remaining-balance commitment
\[
C_{\mathsf{BP}} = m \cdot h_4 + r_{\mathsf{BP}} \cdot h_0.
\]
To prove $s \ge 1$ without disclosing $s$, the client computes an auxiliary commitment to $s - 1$:
\[
C_{s-1} = (s - 1) \cdot h_5 + r_{s-1} \cdot h_0, \quad r_{s-1} \xleftarrow{\$} \mathbb{Z}_{q_{\text{BLS}}}.
\]
It then generates three BEQ proofs:
\begin{itemize}[leftmargin=*]
\item $\pi_{\text{BEQ\_spend}}$ for the refund balance $m$ using BLS generator $h_4$;
\item $\pi_{\text{BEQ\_s}}$ for the hidden spend amount $s$ using BLS generator $h_5$; and
\item $\pi_{\text{BEQ\_s\_minus\_1}}$ for the auxiliary value $s-1$ using BLS generator $h_5$ and commitment $C_{s-1}$.
\end{itemize}
The request discloses $(\mathsf{mode}=\text{``Hidden''}, k_{\mathsf{cur}}, T_{\mathsf{issue}}, \eta, K', C_{\mathsf{BP}}, C_s, C_{s-1})$ and does \emph{not} disclose $s$ or $s_{\mathsf{req}}$.

\paragraph{Zero-Knowledge Proof $\pi_{\mathsf{spend}}$ (Hidden Mode).}
Hidden mode extends the Exact-mode witness set with the scaled hidden amount $\widetilde{a} = s r_1$ and the scaled commitment blinder $x = r_s r_1$. Let $\rho_{\widetilde{a}}, \rho_x \xleftarrow{\$} \mathbb{Z}_q$ be the corresponding blinders and define $\rho_m = \rho_{\tilde{c}} - \rho_{\widetilde{a}}$. The prover computes
\begin{align*}
T_{\mathsf{scale},S} &= \rho_{r_1} \cdot C_s, & T_S &= \rho_{\widetilde{a}} \cdot h_5 + \rho_x \cdot h_0, \\
T_{\mathsf{scale},R} &= \rho_{r_1} \cdot K', & T_{\mathsf{refund}} &= \rho_m \cdot h_4 + \rho_u \cdot h_1 + \rho_{r_1} \cdot T_{\mathsf{issue}} \cdot h_2 + \rho_v \cdot h_0, \\
T_{\mathsf{scale},BP} &= \rho_{r_1} \cdot C_{\mathsf{BP}}, & T_{\mathsf{BP}} &= \rho_m \cdot h_4 + \rho_w \cdot h_0,
\end{align*}
and responses
\begin{align*}
z_{\widetilde{a}} &= \rho_{\widetilde{a}} + c \cdot (s r_1), & z_x &= \rho_x + c \cdot (r_s r_1).
\end{align*}
The prover defines the hidden refund-balance response as
\[
z_m = z_{\tilde{c}} - z_{\widetilde{a}},
\]
which corresponds to $(c_{\mathsf{bal}} - s) r_1 = m r_1$. The Fiat--Shamir challenge is recomputed over the Hidden-mode statement and \textbf{must} bind the mode tag, $C_s$, and $\pi_{\text{BEQ\_s}}$:
\[
c = \mathcal{H}_{\mathsf{scalar}}\left(
\begin{array}{c}
H_{\mathsf{ctx}} \parallel \mathsf{pk}_S \parallel \text{``Hidden''} \parallel k_{\mathsf{cur}} \parallel T_{\mathsf{issue}} \parallel \eta \parallel K' \parallel C_{\mathsf{BP}} \parallel C_s \parallel \pi_{\text{BEQ\_spend}} \parallel \pi_{\text{BEQ\_s}} \parallel \\
A' \parallel \bar{A} \parallel T_{\mathsf{BBS}} \parallel T_{\mathsf{scale},S} \parallel T_S \parallel T_{\mathsf{scale},R} \parallel T_{\mathsf{refund}} \parallel T_{\mathsf{scale},BP} \parallel T_{\mathsf{BP}} \parallel \text{``Spend''}
\end{array}
\right).
\]
On the wire, $z_{\widetilde{a}}$ is serialized as the field
\texttt{z\_amt\_tilde}.

\paragraph{Server Verification (Hidden Mode).}
The server performs the same nonce, double-spend, BBS+, BEQ-for-$m$, and pairing checks as in Exact mode, but with the following mode-specific requirements:
\begin{enumerate}[leftmargin=*]
\item Require $\mathsf{mode} = \text{``Hidden''}$; require $(C_s, C_{s-1}, \pi_{\text{BEQ\_s}}, \pi_{\text{BEQ\_s\_minus\_1}})$ to be present; and reject any disclosed Exact-mode-only fields $(s, s_{\mathsf{req}}, T_{\mathsf{scale},T}, T_{\mathsf{total}})$.
\item Compute the spend idempotency key over the Hidden-mode statement, binding $\text{``Hidden''}$, $C_s$, $C_{s-1}$, $\pi_{\text{BEQ\_s}}$, $\pi_{\text{BEQ\_s\_minus\_1}}$, $\pi_{\mathsf{spend}}$, and $\eta$.
\item Recompute the Fiat--Shamir challenge over the Hidden-mode transcript above.
\item Verify the hidden-amount bridge:
  \[
  z_{\widetilde{a}} \cdot h_5 + z_x \cdot h_0 - T_S \stackrel{?}{=} z_{r_1} \cdot C_s - T_{\mathsf{scale},S}.
  \]
\item Verify the refund and range bridges using $z_m = z_{\tilde{c}} - z_{\widetilde{a}}$:
  \begin{align*}
    z_m \cdot h_4 + z_u \cdot h_1 + z_{r_1} \cdot T_{\mathsf{issue}} \cdot h_2 + z_v \cdot h_0 - T_{\mathsf{refund}} &\stackrel{?}{=} z_{r_1} \cdot K' - T_{\mathsf{scale},R}, \\
    z_m \cdot h_4 + z_w \cdot h_0 - T_{\mathsf{BP}} &\stackrel{?}{=} z_{r_1} \cdot C_{\mathsf{BP}} - T_{\mathsf{scale},BP}.
  \end{align*}
\item Verify $\pi_{\text{BEQ\_s}}$ with respect to $C_s$.  Then enforce $s \ge 1$ via the following two-step check:
  \begin{enumerate}[(a)]
  \item Verify the public commitment relation $C_s = C_{s-1} + h_5$; this publicly proves $s = (s-1)+1$ under the Pedersen binding property.
  \item Verify $\pi_{\text{BEQ\_s\_minus\_1}}$ with respect to $C_{s-1}$; this proves $s-1 \in [0, 2^{32}-1]$, hence $s \in [1, 2^{32}]$.  Combined with the upper-bound check from $\pi_{\text{BEQ\_s}}$ (which proves $s \in [0, 2^{32}-1]$), we obtain $s \in [1, 2^{32}-1]$.
  \end{enumerate}
  \textbf{Reject any request that omits $C_{s-1}$ or $\pi_{\text{BEQ\_s\_minus\_1}}$, or for which $C_s \neq C_{s-1} + h_5$.}  Without this check, a malicious client can spend $s=0$ credits, burning no balance, enabling unlimited free requests and potential denial-of-service or state-probing attacks.
\end{enumerate}
If all checks pass, the server issues the same refund token format as in Exact mode and the client finalizes identically. Hidden mode therefore changes only the statement being proved, not the resulting refund credential.

\subsection{Notation and Preliminaries}

\begin{description}[leftmargin=3.8cm,style=nextline]
\item[$\mathbb{G}_{\text{BLS}}$] BLS12-381 G1 group (written additively).
\item[$\mathbb{G}_{25519}$] Ristretto group over Curve25519 (written additively).
\item[$q_{\text{BLS}}$] Prime order of $\mathbb{G}_{\text{BLS}}$ ($\approx 2^{255}$).
\item[$q_{25519}$] Prime order of $\mathbb{G}_{25519}$ ($\approx 2^{252}$).
\item[$h_{\mathsf{val}}, h_0 \in \mathbb{G}_{\text{BLS}}$] Public BLS12-381 generators for committed value and blinding; instantiate $h_{\mathsf{val}} = h_4$ for balances and $h_{\mathsf{val}} = h_5$ for hidden spend amounts.
\item[$G, H \in \mathbb{G}_{25519}$] Standard Ristretto Pedersen generators.
\item[$\mathcal{H}_{128}$] Hash function $\{0,1\}^* \to \{0,1\}^{128}$ for the cross-curve challenge.
\end{description}

We rely on the \emph{Discrete Logarithm (DL)} assumption in both groups and on the random oracle model for hash functions. Because the two groups are defined over different fields, we assume there is no efficiently computable homomorphism between them; this is guaranteed by the distinct curve parameters.

\subsection{Integer Arithmetic and Bounds}

All Sigma protocol operations involving the secret value $m$ are performed over the integers $\mathbb{Z}$ rather than modulo a group order. We define the following bounds:

\begin{itemize}
\item $m \in [0, 2^{32}-1]$ (enforced by the Bulletproof).
\item The Fiat--Shamir challenge $c$ is a 128-bit integer, i.e.\ $c \in [0, 2^{128}-1]$.
\item The blinding integer $u_m$ is drawn uniformly from $[0, 2^{240}-1]$.
\end{itemize}

The response $z = u_m + c \cdot m$ therefore satisfies $0 \le z < 2^{240} + 2^{128} \cdot (2^{32}-1) < 2^{241}$. Crucially, $2^{241} < q_{25519} < q_{\text{BLS}}$, so $z$ is strictly smaller than the order of both groups. Consequently, interpreting $z$ as a scalar modulo $q_{\text{BLS}}$ or $q_{25519}$ is unambiguous: no modular reduction ever occurs, and the integer value of $z$ is preserved exactly in both groups.

\subsection{The Batched Equality Proof Protocol}
\label{sec:crosscurve}

For a public value generator $h_{\mathsf{val}} \in \{h_4, h_5\}$, the protocol proves knowledge of $(m, r_{\text{BLS}}, r_{25519})$ such that:
\[
\begin{aligned}
C_{\text{BLS}} &= m \cdot h_{\mathsf{val}} + r_{\text{BLS}} \cdot h_0 \quad \in \mathbb{G}_{\text{BLS}}, \\
C_{25519} &= m \cdot G + r_{25519} \cdot H \quad \in \mathbb{G}_{25519},
\end{aligned}
\]
and that $m \in [0, 2^{32}-1]$.

\subsubsection{Prover Algorithm}

\textbf{Input:} $m \in [0, 2^{32}-1]$, a public generator $h_{\mathsf{val}} \in \{h_4,h_5\}$, $r_{\text{BLS}} \in \mathbb{Z}_{q_{\text{BLS}}}$, $r_{25519} \in \mathbb{Z}_{q_{25519}}$, context string $\text{ctx}$, and optional auxiliary commitments $\text{comm}_1,\dots,\text{comm}_k \in \mathbb{G}_{\text{BLS}}$ (e.g., BBS+ proof components).

\begin{enumerate}[leftmargin=*]
\item \textbf{Range Proof:}
  Generate a Dalek Bulletproof $\pi_{\text{range}}$ for the Ristretto commitment $C_{25519}$ proving that $m \in [0, 2^{32}-1]$. Because $\pi_{\text{range}}$ depends only on the statement $C_{25519}$ and not on the Sigma-protocol blinders chosen in step~2, it \textbf{must} be generated before the rejection-sampling loop: the Fiat--Shamir challenge $c$ in step~3 binds $\pi_{\text{range}}$, so $\pi_{\text{range}}$ must be fixed before any blinder is committed to. The Bulletproof transcript is initialized with domain separator \texttt{"EBUT:BEQ:Range"} and includes $\mathsf{len}(\text{ctx}) \parallel \text{ctx}$ followed by $\mathsf{len}(\text{comm}_i) \parallel \text{comm}_i$ for each auxiliary commitment, using the same length-prefixed encoding mandated for the Fiat--Shamir challenge below.

\item \textbf{Commitments (with rejection sampling):}
  Repeat the following until acceptance:
  \begin{enumerate}
  \item Choose random integers
    \[
    u_m \xleftarrow{\$} [0, 2^{240}-1], \qquad
    u_{\text{BLS}} \xleftarrow{\$} \mathbb{Z}_{q_{\text{BLS}}}, \qquad
    u_{25519} \xleftarrow{\$} \mathbb{Z}_{q_{25519}}.
    \]
  \item Compute
    \[
    R_{\text{BLS}} = u_m \cdot h_{\mathsf{val}} + u_{\text{BLS}} \cdot h_0 \in \mathbb{G}_{\text{BLS}}, \qquad
    R_{25519} = u_m \cdot G + u_{25519} \cdot H \in \mathbb{G}_{25519}.
    \]
  \item Compute the challenge $c$ as in step~3 below.
  \item Set $z = u_m + c \cdot m$ (as an integer).
  \item \textbf{Rejection test:} if $z > 2^{240}-1$, \textbf{abort and restart} from the beginning of this loop with fresh randomness.
  \item Otherwise \textbf{accept} and continue.
  \end{enumerate}
  The expected number of iterations is less than~2; rejection occurs with probability at most $(2^{32}-1) \cdot (2^{128}-1) / 2^{240} < 2^{-80}$, where $c_{\max} \le 2^{32}-1$ per Section~4.1 and $c < 2^{128}$.

\item \textbf{Challenge:}
  Compute the 128-bit challenge using length-prefixed encoding for all variable-length fields to prevent byte-shifting transcript collisions (the same requirement mandated for the idempotency cache in Section~7.4):
  \[
  \begin{aligned}
  c = \mathcal{H}_{128}\!\bigl(&\text{EBUT:BEQ:Challenge} \parallel \mathsf{len}(\text{ctx}) \parallel \text{ctx} \parallel h_{\mathsf{val}} \parallel h_0 \parallel G \parallel H \\
  &\parallel C_{\text{BLS}} \parallel C_{25519} \parallel R_{\text{BLS}} \parallel R_{25519} \\
  &\parallel \mathsf{len}(\pi_{\text{range}}) \parallel \pi_{\text{range}} \\
  &\parallel \mathsf{len}(\text{comm}_1) \parallel \text{comm}_1 \parallel \dots \parallel \mathsf{len}(\text{comm}_k) \parallel \text{comm}_k\bigr).
  \end{aligned}
  \]
  where $\mathsf{len}(\cdot)$ denotes a fixed-width 4-byte big-endian encoding of the byte length of the following field. The public generators $h_{\mathsf{val}}, h_0 \in \mathbb{G}_{\text{BLS}}$ (48 bytes each) and $G, H \in \mathbb{G}_{25519}$ (32 bytes each) are fixed size and \textbf{must} be bound into the transcript to prevent parameter-substitution attacks; they do not require length prefixes. The remaining curve group elements $C_{\text{BLS}}$, $C_{25519}$, $R_{\text{BLS}}$, $R_{25519}$ are always fixed size (48, 32, 48, and 32 bytes respectively) and do not require length prefixes; however, $\text{ctx}$, $\pi_{\text{range}}$, and each $\text{comm}_i$ are variable-length and \textbf{must} be length-prefixed. Raw concatenation of variable-length fields without delimiters is \textbf{prohibited}: an attacker can shift bytes between adjacent variable-length fields to produce an identical transcript hash for two distinct proof statements, breaking binding.

\item \textbf{Responses:}
  Use the $z = u_m + c \cdot m$ accepted in step~2 (over the integers, not reduced modulo any group order).
  Compute modular responses:
  \[
  z_{\text{BLS}} = u_{\text{BLS}} + c \cdot r_{\text{BLS}} \pmod{q_{\text{BLS}}}, \qquad
  z_{25519} = u_{25519} + c \cdot r_{25519} \pmod{q_{25519}}.
  \]
  Here $c$ is interpreted as an element of the respective scalar fields by reducing modulo $q_{\text{BLS}}$ and $q_{25519}$ respectively; because $c < 2^{128} \ll \min(q_{\text{BLS}}, q_{25519})$, this reduction is the identity map.

\item \textbf{Output:}
  The canonical proof is
  \[
  \pi_{\text{BEQ}} = (C_{25519}, R_{\text{BLS}}, R_{25519}, z, z_{\text{BLS}}, z_{25519}, \pi_{\text{range}}).
  \]

  \textbf{Transcript compression (reference implementation):}
  The reference implementation omits $R_{\text{BLS}}$ and $R_{25519}$ from the wire format, transmitting instead only the challenge tuple
  \[
  (c_{128}, C_{25519}, z, z_{\text{BLS}}, z_{25519}, \pi_{\text{range}}).
  \]
  The verifier reconstructs:
  \[
  R_{\text{BLS}} = z \cdot h_{\mathsf{val}} + z_{\text{BLS}} \cdot h_0 - c \cdot C_{\text{BLS}}, \quad
  R_{25519} = z \cdot G + z_{25519} \cdot H - c \cdot C_{25519},
  \]
  and re-hashes them to verify the challenge. This saves 80 bytes (48 bytes for $R_{\text{BLS}}$ and 32 bytes for $R_{25519}$) per proof without any loss of security.
\end{enumerate}

\subsubsection{Verifier Algorithm}

\textbf{Input:} $\pi_{\text{BEQ}}$, $C_{\text{BLS}}$, context $\text{ctx}$, auxiliary commitments $\text{comm}_1,\dots,\text{comm}_k$.

\begin{enumerate}[leftmargin=*]
\item \textbf{Bound Check:}
  \textbf{Deserialize $z$ as a raw unsigned big-endian integer occupying exactly 31 bytes; do not pass it to any curve-scalar deserialization routine.}  Many elliptic-curve libraries (e.g., Dalek's \texttt{Scalar::from\_bytes\_mod\_order}, blstrs' scalar deserializer) silently reduce their input modulo the group order on deserialization.  Calling such a routine on $z$ before performing the bounds check will silently map any adversarially crafted $z \ge q_{25519}$ to a value in $[0, q_{25519})$, causing the Ristretto verification equation to accept with a different effective integer than the BLS equation uses, and breaking the cross-curve equality binding.  The bounds check must therefore be the \emph{first} operation performed on the raw bytes:
  Verify that $0 \le z \le 2^{240}-1$ (as an integer). If not, reject. This tighter bound (compared to checking only $z < q_{25519}$, where $q_{25519} \approx 2^{252}$) is enforced by the prover's rejection-sampling step and is required for the zero-knowledge guarantee: without it, a large $z > 2^{240}-1$ would reveal that $c \cdot m$ is large, leaking a lower bound on the secret $m$.  Only after the bounds check passes should the implementation derive the field-reduced values $z_{\text{BLS}} = z \bmod q_{\text{BLS}}$ and $z_{25519} = z \bmod q_{25519}$ for use in the scalar-multiplication equations below.

\item \textbf{Recompute Challenge:}
  Compute $c$ using the identical length-prefixed encoding as the prover:
  \[
  \begin{aligned}
  c = \mathcal{H}_{128}\!\bigl(&\text{EBUT:BEQ:Challenge} \parallel \mathsf{len}(\text{ctx}) \parallel \text{ctx} \parallel h_{\mathsf{val}} \parallel h_0 \parallel G \parallel H \\
  &\parallel C_{\text{BLS}} \parallel C_{25519} \parallel R_{\text{BLS}} \parallel R_{25519} \\
  &\parallel \mathsf{len}(\pi_{\text{range}}) \parallel \pi_{\text{range}} \\
  &\parallel \mathsf{len}(\text{comm}_1) \parallel \text{comm}_1 \parallel \dots \parallel \mathsf{len}(\text{comm}_k) \parallel \text{comm}_k\bigr).
  \end{aligned}
  \]

\item \textbf{Verify BLS12-381 Equation:}
  Let $z_{\text{BLS,red}} = z \bmod q_{\text{BLS}}$ (derived from the already-bounds-checked integer $z$).
  \[
  z_{\text{BLS,red}} \cdot h_{\mathsf{val}} + z_{\text{BLS}} \cdot h_0 \stackrel{?}{=} R_{\text{BLS}} + c \cdot C_{\text{BLS}} \quad \text{in } \mathbb{G}_{\text{BLS}}.
  \]

\item \textbf{Verify Ristretto Equation:}
  Let $z_{25519,\text{red}} = z \bmod q_{25519}$ (derived from the already-bounds-checked integer $z$; since $z \le 2^{240} < q_{25519}$, this reduction is the identity).
  \[
  z_{25519,\text{red}} \cdot G + z_{25519} \cdot H \stackrel{?}{=} R_{25519} + c \cdot C_{25519} \quad \text{in } \mathbb{G}_{25519}.
  \]

\item \textbf{Verify Bulletproof:}
  Verify $\pi_{\text{range}}$ with respect to $C_{25519}$ using the same transcript initialization as the prover.
\end{enumerate}

If all checks pass, the verifier accepts.

\subsection{Security Analysis of BEQ}

\begin{theorem}[Completeness]
If the prover follows the protocol honestly, the verifier always accepts.
\end{theorem}
\begin{proof}
The rejection-sampling loop guarantees that the accepted $z = u_m + c \cdot m$ satisfies $0 \le z \le 2^{240}-1 < q_{25519}$ (where $q_{25519} = 2^{252} + 27742317777372353535851937790883648493$), so the verifier's bound check passes. Substituting the definitions:
\[
z \cdot h_{\mathsf{val}} + z_{\text{BLS}} \cdot h_0 = (u_m + c m) h_{\mathsf{val}} + (u_{\text{BLS}} + c r_{\text{BLS}}) h_0 = R_{\text{BLS}} + c \cdot C_{\text{BLS}},
\]
and similarly for the Ristretto equation. The Bulletproof is complete by the underlying Dalek implementation. Thus, all checks succeed.
\end{proof}

\begin{theorem}[Knowledge Soundness]
Assume the discrete logarithm problem is hard in both $\mathbb{G}_{\text{BLS}}$ and $\mathbb{G}_{25519}$, and that the Bulletproof is computationally sound. Let $\mathcal{A}$ be a PPT adversary making at most $q_{\text{RO}}$ queries to $\mathcal{H}_{128}$ and producing a valid proof $\pi_{\text{BEQ}}$ with probability $\epsilon$. Then there exists an extractor $\mathcal{E}$ that, given black-box access to $\mathcal{A}$, outputs a witness $(m, r_{\text{BLS}}, r_{25519})$ such that $m \in [0,2^{32}-1]$ and the commitment equations hold, with success probability at least
\[
\frac{\epsilon^2}{q_{\text{RO}}} - \frac{\epsilon}{2^{128}} - \advantage{DL}(\mathcal{B}) - \advantage{BP\text{-}sound}(\mathcal{B}').
\]
\end{theorem}
\begin{proof}
Consider the interactive version of the protocol where the challenge $c$ is chosen uniformly from $\{0,1\}^{128}$ by the verifier. By the general forking lemma~\cite{PointchevalStern00} with challenge space size $|\mathcal{C}| = 2^{128}$, if $\mathcal{A}$ succeeds with probability $\epsilon$, then with probability at least $\epsilon^2/q_{\text{RO}} - \epsilon/2^{128}$ we can rewind $\mathcal{A}$ to obtain two accepting transcripts with the same first message $(C_{25519}, R_{\text{BLS}}, R_{25519}, \pi_{\text{range}})$ and distinct challenges $c \neq c'$. Let the responses be $(z, z_{\text{BLS}}, z_{25519})$ and $(z', z_{\text{BLS}}', z_{25519}')$. From the verification equations we have:
\begin{align*}
(z - z') \cdot h_{\mathsf{val}} + (z_{\text{BLS}} - z_{\text{BLS}}') \cdot h_0 &= (c - c') \cdot C_{\text{BLS}} \quad \text{in } \mathbb{G}_{\text{BLS}}, \\
(z - z') \cdot G + (z_{25519} - z_{25519}') \cdot H &= (c - c') \cdot C_{25519} \quad \text{in } \mathbb{G}_{25519}.
\end{align*}
Define $\Delta z = z - z'$, $\Delta z_{\text{BLS}} = z_{\text{BLS}} - z_{\text{BLS}}'$, $\Delta z_{25519} = z_{25519} - z_{25519}'$, and $\Delta c = c - c'$. Since $c \neq c'$, $\Delta c \neq 0$ and $|\Delta c| < 2^{128}$. The verifier's bound check ensures $0 \le z, z' < q_{25519}$, so $|\Delta z| < q_{25519}$. Because $q_{25519} \approx 2^{252}$, the integer division $\Delta z / \Delta c$ yields a unique integer $m$ with $|m| < q_{25519}$. Moreover, since $z \ge 0$ and $c \ge 0$, we can show $m \ge 0$ by considering the two cases of $\Delta c$ sign.

Now consider the first equation in $\mathbb{G}_{\text{BLS}}$. Since $|\Delta z| < q_{\text{BLS}}$, the equality holds over integers when interpreted as scalars modulo $q_{\text{BLS}}$. Thus we can write
\[
C_{\text{BLS}} = \left(\frac{\Delta z}{\Delta c}\right) \cdot h_{\mathsf{val}} + \left(\frac{\Delta z_{\text{BLS}}}{\Delta c}\right) \cdot h_0 \pmod{q_{\text{BLS}}}.
\]
Set $m = \frac{\Delta z}{\Delta c}$ (computed over integers) and $r_{\text{BLS}} = \frac{\Delta z_{\text{BLS}}}{\Delta c} \bmod q_{\text{BLS}}$. Similarly, from the second equation we obtain $m$ and $r_{25519} = \frac{\Delta z_{25519}}{\Delta c} \bmod q_{25519}$. The Bulletproof soundness guarantees $m \in [0,2^{32}-1]$ except with probability $\advantage{BP\text{-}sound}(\mathcal{B}')$. If the adversary produces proofs that verify but extraction fails, we can embed a DL challenge into the generators and extract a discrete logarithm from the discrepancy; the reduction is standard and omitted. The bound follows.
\end{proof}

\begin{theorem}[Zero-Knowledge]
Under the random oracle model, the non-interactive BEQ protocol with rejection sampling is statistically zero-knowledge with distinguishing advantage bounded by $2^{-80}$. Moreover, the support of $z$ in real executions is exactly $[0, 2^{240}-1]$---the same as the simulator---so no transcript component is outside the simulator's range.
\end{theorem}
\begin{proof}
We construct a simulator $\mathcal{S}$ that, on input $(C_{\text{BLS}}, \text{ctx})$ and without knowledge of the witness, produces a transcript $\pi_{\text{BEQ}} = (C_{25519}, R_{\text{BLS}}, R_{25519}, z, z_{\text{BLS}}, z_{25519}, \pi_{\text{range}})$ as follows:
\begin{enumerate}
\item Choose a random challenge $c \xleftarrow{\$} \{0,1\}^{128}$.
\item Choose random $z \xleftarrow{\$} [0, 2^{240}-1]$, $z_{\text{BLS}} \xleftarrow{\$} \mathbb{Z}_{q_{\text{BLS}}}$, $z_{25519} \xleftarrow{\$} \mathbb{Z}_{q_{25519}}$.
\item Compute $R_{\text{BLS}} = z \cdot h_{\mathsf{val}} + z_{\text{BLS}} \cdot h_0 - c \cdot C_{\text{BLS}}$ and $R_{25519} = z \cdot G + z_{25519} \cdot H - c \cdot C_{25519}$, where $C_{25519}$ is an arbitrary commitment to $0$.
\item Generate a simulated Bulletproof $\pi_{\text{range}}$ for $C_{25519}$ with range $[0,2^{32}-1]$ using the honest-verifier zero-knowledge property of Dalek Bulletproofs (by programming the internal random oracle).
\item Program $\mathcal{H}_{128}$ to return $c$ on the appropriate input.
\end{enumerate}
In a real execution with rejection sampling, the prover discards any $u_m$ for which $z = u_m + c \cdot m > 2^{240}-1$. For any fixed challenge $c$ and message $m$, the map $u_m \mapsto z = u_m + c \cdot m$ is a bijection on integers, so conditioning on $z \le 2^{240}-1$ is equivalent to conditioning on $u_m \le 2^{240}-1 - c \cdot m$. Since $u_m$ is uniform over $[0, 2^{240}-1]$ and the accepted values fill a contiguous sub-interval of length $2^{240} - c \cdot m$, the accepted $z$ is uniform over $[0, 2^{240}-1 - c \cdot m] \subseteq [0, 2^{240}-1]$. The simulated $z$ is uniform over the full $[0, 2^{240}-1]$. The statistical distance between these distributions is $c \cdot m / 2^{240} \le (2^{128}-1)(2^{32}-1)/2^{240} < 2^{-80}$, which is negligible---but more importantly, both distributions are supported on $[0, 2^{240}-1]$, so no single value of $z$ is impossible in simulation yet possible in real execution or vice versa. The other components are identically distributed. Hence, the simulation achieves negligible statistical distance, and the distinguishing advantage is at most $2^{-80}$.
\end{proof}

\subsubsection{Transcript Binding}
To prevent splicing attacks where a Bulletproof generated for one context is reused in another, the prover \textbf{must} include the BBS+ commitments ($A'$, $\bar{A}$, $T_{\mathsf{BBS}}$) and any bridge commitments in the Bulletproof transcript. Concretely, the transcript is initialized as:
\begin{verbatim}
Transcript::new(b"EBUT:BEQ:Range");
transcript.append_message(b"ctx", context);
for comm in auxiliary_commitments {
    transcript.append_message(b"comm", comm.to_compressed());
}
\end{verbatim}
This binds the range proof to the specific BBS+ proof instance, as required by the EBUT construction.

\subsection{Integration into EBUT}
The cross-curve batched equality proof replaces all standalone range proofs in the EBUT protocol.
\begin{itemize}
\item \textbf{Epoch Refresh:} The commitment $C_{\Delta}$ to the expiry delta $\Delta = E_{\mathsf{max}} - T$ is bound to a Ristretto commitment and a range proof via $\pi_{\text{BEQ\_exp}}$. The auxiliary commitments include $A', \bar{A}, T_{\mathsf{BBS}}, N_T, K_{\mathsf{daily}}$.
\item \textbf{Spend Phase (Exact mode):} The commitment $C_{\mathsf{BP}}$ to the remaining balance $m = c_{\mathsf{bal}} - s$ is linked to a Ristretto commitment via $\pi_{\text{BEQ\_spend}}$ instantiated with $h_{\mathsf{val}} = h_4$. Auxiliary commitments include $A', \bar{A}, T_{\mathsf{BBS}}, K', C_{\mathsf{total}}$.
\item \textbf{Spend Phase (Hidden mode):} The refund-balance proof again uses $h_{\mathsf{val}} = h_4$, while the hidden amount proof $\pi_{\text{BEQ\_s}}$ uses $h_{\mathsf{val}} = h_5$ for the separate amount commitment $C_s$.
\end{itemize}
The proof size is approximately 672 bytes for the Bulletproof (single 32-bit range) plus the Sigma protocol responses ($\approx 160$ bytes), totaling under 1~KB. Verification time is dominated by the Bulletproof ($\approx 1$--$2$~ms) and the two multi-scalar multiplications, which are negligible by comparison.

\section{Security Analysis of EBUT}
\label{sec:security-ebut}

We now prove that the EBUT construction satisfies the three security properties defined in Section~\ref{sec:defs}: unforgeability, unlinkability, and rate-limit soundness. All proofs are in the random oracle model and rely on the hardness assumptions stated in Section~2. We assume the cross-curve BEQ is used for range proofs; its security properties have been established in Section~\ref{sec:crosscurve}.

\subsection{Unforgeability}

\begin{theorem}[Unforgeability]
Let $\mathcal{A}$ be a PPT adversary against the $\mathsf{UNF}$ game making at most $q_I$ issuance queries, $q_R$ refresh queries, $q_S$ spend queries, and $q_{\text{RO}}$ random oracle queries. Then there exist PPT adversaries $\mathcal{B}_{\text{SDH}}$ and $\mathcal{B}_{\text{DL}}$ such that
\begin{align*}
\advantage{UNF}(\mathcal{A}) \le \advantage{q\text{-}SDH}(\mathcal{B}_{\text{SDH}}) + \advantage{DL}(\mathcal{B}_{\text{DL}}) + \frac{(q_I+q_R+q_S+q_{\text{RO}})^2}{2q} \\
+ \frac{q_R+q_S}{q} + \frac{1}{2^{128}} + \negl(\lambda),
\end{align*}
where $q$ is the group order.
\end{theorem}
\begin{proof}
We proceed through a sequence of games $\mathsf{G}_0,\dots,\mathsf{G}_8$. Let $\Pr[S_i]$ denote the probability that $\mathcal{A}$ wins in Game $\mathsf{G}_i$.

\paragraph{Game $\mathsf{G}_0$:} The original $\mathsf{UNF}$ game. The challenger $\mathcal{C}$ runs $\mathsf{Setup}$, generates $(\mathsf{sk}_S,\mathsf{pk}_S)=(y,g_2^y)$, and answers all oracles honestly.

\paragraph{Game $\mathsf{G}_1$ (Simulation of zero-knowledge proofs):} $\mathcal{C}$ replaces all zero-knowledge proofs ($\pi_{\mathsf{refresh}}$, $\pi_{\mathsf{spend}}$, and the internal BEQ proofs) with simulated proofs using the simulators from the ZK theorems. The simulation is perfect under DL and the statistical ZK of BEQ. Any distinguisher gives an adversary against DL or the statistical test. Thus,
\[
|\Pr[S_1] - \Pr[S_0]| \le \advantage{DL}(\mathcal{B}_1) + \frac{q_{\text{RO}}^2}{2q} + 2^{-80}.
\]

\paragraph{Game $\mathsf{G}_2$ (No random oracle collisions):} $\mathcal{C}$ aborts if a collision occurs in the outputs of $\mathcal{H}_{\mathsf{scalar}}$ or $\mathcal{H}_{128}$. The probability is at most $q_{\text{RO}}^2/2q + q_{\text{RO}}^2/2^{129}$. Hence,
\[
|\Pr[S_2] - \Pr[S_1]| \le \frac{q_{\text{RO}}^2}{2q} + \frac{q_{\text{RO}}^2}{2^{129}}.
\]

\paragraph{Game $\mathsf{G}_3$ (Extraction from refresh proofs):} For each valid refresh request, $\mathcal{C}$ applies the knowledge extractor for the BBS+ proof and the Schnorr equalities. By the forking lemma on the interactive version (challenge space $\mathbb{Z}_q$ for the BBS+ proof, and the Schnorr part), with probability at least $\epsilon^2/q_{\text{RO}} - \epsilon/q$, $\mathcal{C}$ extracts the master secret $k_{\mathsf{sub}}$ and the BBS+ witness. If extraction fails, the proof is a forgery, which reduces to DL or $q$-SDH. The soundness error adds at most $\advantage{DL}(\mathcal{B}_2) + q_R/q$. Thus,
\[
|\Pr[S_3] - \Pr[S_2]| \le \advantage{DL}(\mathcal{B}_2) + \frac{q_R}{q}.
\]

\paragraph{Game $\mathsf{G}_4$ (Binding nullifiers to $k_{\mathsf{sub}}$):} $\mathcal{C}$ enforces the invariant that for any two accepted refresh requests with the same $\mathsf{pk}_C$ and epoch $T$, the extracted $k_{\mathsf{sub}}$ are equal. If not, we show a reduction to DL. Suppose $N_T = k \cdot H_T$ and $N_T' = k' \cdot H_T$ with $k \neq k'$. Then $(k - k') \cdot H_T = N_T - N_T'$. Given a DL challenge $H_T = g_1^a$, we can embed it and extract $a$ from the known $k,k',N_T,N_T'$. The probability of this event is bounded by $\advantage{DL}(\mathcal{B}_3)$. Thus,
\[
|\Pr[S_4] - \Pr[S_3]| \le \advantage{DL}(\mathcal{B}_3).
\]

\paragraph{Game $\mathsf{G}_5$ (Simulation of issuance via BBS+ signing oracle):} $\mathcal{C}$ replaces the honest issuance in $\mathcal{O}_{\mathsf{Issue}}$ with queries to a BBS+ signing oracle $\mathcal{O}_{\mathsf{Sign}}$. When $\mathcal{A}$ requests issuance for $\mathsf{pk}_C$, $\mathcal{C}$ extracts $k_{\mathsf{sub}}$ from the proof of knowledge (standard Schnorr) and forwards the blinded commitment to $\mathcal{O}_{\mathsf{Sign}}$, which returns a signature on $(k_{\mathsf{sub}}, c_{\mathsf{max}}, E_{\mathsf{max}})$ under the secret key $y$ (unknown to $\mathcal{C}$). The simulation is perfect, so $\Pr[S_5] = \Pr[S_4]$.

\paragraph{Game $\mathsf{G}_6$ (Detection of master token forgeries):} $\mathcal{C}$ checks if $\mathcal{A}$ outputs a valid master token for a $\mathsf{pk}_C^*$ that was never issued, or more tokens than $q_{\mathsf{issue}}$. This directly breaks EUF-CMA of BBS+. We construct $\mathcal{B}_{\text{SDH}}$ that simulates $\mathsf{G}_6$ and outputs the forgery. Hence,
\[
\Pr[\text{Case A win in }\mathsf{G}_6] \le \advantage{q\text{-}SDH}(\mathcal{B}_{\text{SDH}}).
\]

\paragraph{Game $\mathsf{G}_7$ (Multiple daily tokens):} $\mathcal{C}$ checks if $\mathcal{A}$ produced more than one daily token for the same $\mathsf{pk}_C^*$ and epoch $T$. By $\mathsf{G}_4$, all such tokens have the same $k_{\mathsf{sub}}$ and thus identical nullifier $N_T = k_{\mathsf{sub}} \cdot H_T$. The server's $\DB_{\mathsf{epoch}}$ prevents duplicate nullifiers. To win, $\mathcal{A}$ must forge the freshness check, which is impossible in the simulation. The probability is bounded by the soundness error:
\[
\Pr[\text{Case B win in }\mathsf{G}_7] \le \frac{q_R}{q} + \negl(\lambda).
\]

\begin{remark}[Ideal-DB assumption and practical partial-write safety]
The proof of $\mathsf{G}_7$ assumes that $\DB_{\mathsf{epoch}}$ is updated \emph{atomically and instantaneously}: in the ideal model there are no crashes, no network partitions, and no partial writes between the idempotency cache and the nullifier store.  A real implementation must bridge this gap.  The key mechanism is to derive the Refresh idempotency cache key from $N_T$ rather than from the proof transcript $\pi_{\mathsf{refresh}}$ (Section~\ref{sec:idempotency}): because $N_T$ is cryptographically bound to the client's master key and the epoch, an adversary cannot alter it between retries.  Consequently, even if the server crashes after writing the cache entry but before inserting $N_T$ into $\DB_{\mathsf{epoch}}$, any subsequent request carrying the same $N_T$---regardless of any change in $K_{\mathsf{daily}}$ or $r_{\mathsf{daily}}$---will unconditionally hit the cache and return the original response.  The formal bound above is therefore achievable in practice when the implementation obeys the write-ordering invariant described in Section~\ref{sec:idempotency}.
\end{remark}

\paragraph{Game $\mathsf{G}_8$ (Overspending):} Suppose $\mathcal{A}$ attempts to overspend. Each spend proof includes a BEQ proof that the new balance $m \ge 0$ and Schnorr equalities ensuring $m = c_{\mathsf{bal}} - s$. By induction on the number of spends from a single daily token, the total spent cannot exceed $c_{\mathsf{max}}$. To overspend, $\mathcal{A}$ must either (a) double-spend a token nullifier $k_{\mathsf{cur}}$ (prevented by $\DB_{\mathsf{spend}}$), (b) forge a spend proof (breaks DL), or (c) spend from multiple daily tokens derived from the same master token in the same epoch, which is Case B and already bounded. Thus,
\[
\Pr[\text{Case C win in }\mathsf{G}_8] \le \Pr[\text{Case B win}] + \advantage{DL}(\mathcal{B}_4).
\]

Combining the games, we set $\mathcal{B}_{\text{DL}}$ to be the adversary with maximum advantage among $\mathcal{B}_1,\dots,\mathcal{B}_4$. The total advantage is bounded as stated.
\end{proof}

\subsection{Unlinkability}
\label{sec:unlinkability}

\begin{theorem}[Unlinkability]
\label{thm:unlinkability}
Let $\mathcal{A} = (\mathcal{A}_1, \mathcal{A}_2)$ be a PPT adversary against the $\mathsf{UNLINK}$ game making at most $q_{\text{RO}}$ random oracle queries, $q_{\mathsf{refresh}}$ challenge refresh queries, and $q_{\mathsf{spend}}$ challenge spend queries. Then there exists an adversary $\mathcal{B}_{\text{SXDH}}$ such that
\[
\advantage{UNLINK}(\mathcal{A}) \le 2 \cdot \advantage{SXDH}(\mathcal{B}_{\text{SXDH}}) + \frac{q_{\text{RO}}^2}{q} + \frac{q_{\mathsf{refresh}}+q_{\mathsf{spend}}}{q} + \frac{1}{2^{128}} + \negl(\lambda).
\]
\end{theorem}
\begin{proof}
We proceed via a sequence of games $\mathsf{H}_0,\dots,\mathsf{H}_4$. Let $\Pr[S_{\mathsf{H}_i}]$ denote the probability that $\mathcal{A}$ outputs the correct bit $b$.

\paragraph{Game $\mathsf{H}_0$:} The real UNLINK game.

\paragraph{Game $\mathsf{H}_1$ (Simulation of ZK proofs):} $\mathcal{C}$ simulates all ZK proofs in the challenge phase. The difference is bounded by the ZK error:
\[
|\Pr[S_{\mathsf{H}_1}] - \Pr[S_{\mathsf{H}_0}]| \le \advantage{DL}(\mathcal{B}_{\text{DL}}) + \frac{q_{\text{RO}}^2}{2q} + 2^{-80}.
\]

\paragraph{Game $\mathsf{H}_2$ (Simulator Trapdoor \& Nullifier Randomization):}
In this game the adversary $\mathcal{A}$ (acting as a malicious server) outputs an arbitrary public key $\mathsf{pk}_S$ together with a valid proof $\pi_S = (T_y, z_y)$ as required by the extended $\mathsf{Server.KeyGen}$ (Section~\ref{sec:keygen}).

Because $\pi_S$ is a Fiat--Shamir compiled Schnorr proof in the Random Oracle Model, the Simulator $\mathcal{C}$ applies the standard rewindable extractor on $\mathcal{H}_{\mathsf{scalar}}$ to extract the underlying secret key $y$ satisfying $\mathsf{pk}_S = g_2^y$.  If extraction fails (with probability bounded by the knowledge-soundness error of the Schnorr protocol), $\mathcal{C}$ aborts.

During the Challenge phase, $\mathcal{C}$ replaces $N_T = k_{\mathsf{sub}}^{(b)} \cdot H_T$ with a uniformly random $N_T \xleftarrow{\$} \mathbb{G}_1$.  To simulate a valid BBS+ zero-knowledge transcript without knowing the client's master secret, $\mathcal{C}$ proceeds as follows:
\begin{enumerate}[leftmargin=*,nosep]
  \item $\mathcal{C}$ chooses a random $A' \xleftarrow{\$} \mathbb{G}_1$ and a random challenge $c \xleftarrow{\$} \mathbb{Z}_q$.
  \item $\mathcal{C}$ sets $\bar{A} = A'^y \in \mathbb{G}_1$.  Because $\mathcal{C}$ knows $y$, the pairing equation $e(A', \mathsf{pk}_S) = e(\bar{A}, g_2)$ is perfectly satisfied.
  \item $\mathcal{C}$ computes the forward commitment $T_{\mathsf{BBS}}$ backward from random scalar responses (e.g., $T_{\mathsf{BBS}} = A'^{-z_e} \cdot g_1^{z_{r_1}} \cdot \bar{A}^{-c} \cdots$ for fresh random $z_e, z_{r_1},\dots$).
  \item $\mathcal{C}$ programs the Random Oracle so that the Fiat--Shamir challenge evaluates to $c$.
\end{enumerate}
Because $\mathcal{C}$ extracted the trapdoor $y$, the simulation of the BBS+ transcript is perfectly indistinguishable from a real proof.  The adversary's view therefore depends only on $N_T$.  By the SXDH assumption in $\mathbb{G}_1$, distinguishing $N_T = k_{\mathsf{sub}}^{(b)} \cdot H_T$ from a uniformly random group element is as hard as DDH in $\mathbb{G}_1$; equivalently, any distinguisher yields an SXDH adversary.  Hence,
\[
|\Pr[S_{\mathsf{H}_2}] - \Pr[S_{\mathsf{H}_1}]| \le \advantage{SXDH}(\mathcal{B}_{\text{SXDH}}).
\]

\paragraph{Game $\mathsf{H}_3$ (Randomizing token commitments in spend):} For challenge spend queries, $\mathcal{C}$ replaces the refund commitment $K'$ with a commitment to a random balance (subject to the correct total spent) and simulates the spend proof. The BBS+ randomization ensures that $A'$ and $\bar{A}$ are uniformly distributed among valid signatures, independent of the original token. The Pedersen commitment is perfectly hiding. In hidden-amount mode, the server's view additionally contains only a random commitment $C_s$ and a zero-knowledge proof for the amount relation, which do not leak $s$. Therefore, the view is statistically independent of $b$, and the advantage bound remains unchanged. The change is information-theoretic.

\paragraph{Game $\mathsf{H}_4$ (Ideal game):} $\mathcal{C}$ now ignores $b$ entirely. The view of $\mathcal{A}$ is independent of $b$, so $\Pr[S_{\mathsf{H}_4}] = 1/2$.

Since $\mathsf{H}_3$ and $\mathsf{H}_4$ are identical up to statistical distance, $\Pr[S_{\mathsf{H}_3}] = 1/2$. Summing the bounds yields the theorem.
\end{proof}

\subsection{Rate-Limit Soundness}

\begin{theorem}[Rate-Limit Soundness]
The EBUT construction satisfies rate-limit soundness under the Discrete Logarithm assumption. Concretely, for any PPT adversary $\mathcal{A}$,
\[
\advantage{RLS}(\mathcal{A}) \le \advantage{DL}(\mathcal{B}_{\text{DL}}) + \frac{q_{\text{RO}}^2}{2q} + \frac{q_{\mathsf{refresh}}}{q} + \frac{1}{2^{128}} + \negl(\lambda).
\]
\end{theorem}
\begin{proof}
We analyze the two winning conditions of Game~3.

\paragraph{Case 1: Multiple daily tokens from the same master secret.}
Suppose $\mathcal{A}$ outputs $n \ge 2$ valid refresh tuples for epoch $T$ with distinct nullifiers $N_T^{(1)} \neq N_T^{(2)}$, both verifying under the same $\mathsf{pk}_S$.

By applying the Forking Lemma to the Schnorr equality equations within the BEQ cross-curve proof (with challenge space $2^{128}$), the extractor $\mathcal{E}$ successfully extracts the witnesses $(k_{\mathsf{sub}}^{(1)}, r_1^{(1)})$ and $(k_{\mathsf{sub}}^{(2)}, r_1^{(2)})$ from the two proofs.  If extraction fails, $\mathcal{A}$ has forged the knowledge proof, which reduces to the Discrete Logarithm problem in $\mathbb{G}_1$.  Extraction succeeds with probability at least $\epsilon^2/q_{\text{RO}} - \epsilon/2^{128}$, where $\epsilon$ is the success probability of $\mathcal{A}$.

The verification equations algebraically enforce that $N_T^{(i)} = k_{\mathsf{sub}}^{(i)} \cdot H_T$ for $i \in \{1,2\}$.  Now assume both tokens were derived from the same master secret, i.e., $k_{\mathsf{sub}}^{(1)} = k_{\mathsf{sub}}^{(2)} = k$.  Then
\[
N_T^{(1)} = k \cdot H_T = N_T^{(2)},
\]
which directly contradicts the premise $N_T^{(1)} \neq N_T^{(2)}$.  Therefore $\mathcal{A}$ must have either:
\begin{enumerate}[leftmargin=*,nosep]
  \item bypassed the server's strict $\DB_{\mathsf{epoch}}$ distinctness check (impossible in the ideal model with a correct server implementation), or
  \item forged the zero-knowledge proof of knowledge of $N_T = k \cdot H_T$, which breaks the Discrete Logarithm assumption in $\mathbb{G}_1$.
\end{enumerate}
Thus
\[
\advantage{RLS\text{-Case1}}(\mathcal{A}) \le \advantage{DL}(\mathcal{B}_{\text{DL}}^{(1)}).
\]

\paragraph{Case 2: Overspending.}
Suppose $\mathcal{A}$ produces a set of spend transactions within epoch $T$ that collectively consume more than $c_{\mathsf{max}}$ credits from a single master token without double-spending the same nullifier. Each spend transaction includes a spend proof that reduces to a daily token with balance $c_{\mathsf{max}}$. The spend proof's BEQ range proof ensures $m \ge 0$, and the Schnorr equalities enforce $m = c_{\mathsf{bal}} - s$. By induction on the number of spends, the total spent cannot exceed $c_{\mathsf{max}}$. To overspend, $\mathcal{A}$ must either (a) double-spend a token nullifier (prevented by $\DB_{\mathsf{spend}}$; any collision would imply a forgery of the spend proof), (b) produce a spend proof with an invalid balance (breaks soundness, reduces to DL), or (c) spend from multiple daily tokens derived from the same master secret in the same epoch, which is Case 1 and already bounded.

Thus, the overall advantage is bounded as stated.
\end{proof}

\section{Implementation Guidance}

We now provide detailed, production-oriented guidance for implementing the EBUT protocol. These recommendations are essential for security, interoperability, and performance.

\subsection{BBS+ Library}

Implementers should use a CFRG-compliant BBS+ library that supports blind issuance and zero-knowledge proofs of possession. The proof described in Protocol~1 of the ACT paper uses multiplicative randomization ($A' = A^{r_1}$) and is proven secure under the assumptions of this paper. Ensure that the verification equations match exactly those specified in Section~\ref{sec:construction}. In particular, the pairing check must be $e(A', \mathsf{pk}_S) \stackrel{?}{=} e(\bar{A}, g_2)$ after proper randomization.

\subsection{Cross-Curve Batched Equality Proof Implementation}

The EBUT protocol uses the BEQ to prove that a value $m$ is in $[0,2^{32}-1]$ and is consistent with the BLS12-381 commitment used in the BBS+ proof. To implement this correctly:

\begin{enumerate}[leftmargin=*]
\item \textbf{BLS12-381 Commitment:} Compute $C_{\text{BLS}}$ using the public value generator $h_{\mathsf{val}}$ and the blinding generator $h_0$. Use $h_{\mathsf{val}} = h_4$ for balance commitments and $h_{\mathsf{val}} = h_5$ for hidden spend-amount commitments.
\item \textbf{Ristretto Commitment:} Independently, using the \nolinkurl{curve25519-dalek} library with standard Ristretto generators \nolinkurl{RISTRETTO_BASEPOINT_POINT} (for $G$) and \nolinkurl{RISTRETTO_BASEPOINT_COMPRESSED} (for $H$, derived via hash-to-curve with domain separation), compute $C_{25519} = m \cdot G + r_{25519} \cdot H$ where $r_{25519} \xleftarrow{\$} \mathbb{Z}_{q_{25519}}$.
\item \textbf{BEQ Proof Generation:} Execute the BEQ prover algorithm (Section~\ref{sec:crosscurve}) with inputs $(m, h_{\mathsf{val}}, r_{\text{BLS}}, r_{25519})$, context $\text{ctx} = H_{\text{ctx}}$, and auxiliary commitments being the BBS+ proof components ($A', \bar{A}, T_{\mathsf{BBS}}$, etc.). The prover outputs the tuple $\pi_{\text{BEQ}}$ with components $(C_{25519}, R_{\text{BLS}}, R_{25519}, z, z_{\text{BLS}}, z_{25519}, \pi_{\text{range}})$.
\item \textbf{Verification:} The server verifies $\pi_{\text{BEQ}}$ against $C_{\text{BLS}}$ and the auxiliary commitments. The Dalek Bulletproof verification uses the library's \textbf{default generators}.
\end{enumerate}

\textbf{Critical:} Do \textbf{not} attempt to override the Dalek Bulletproof generators with BLS12-381 points. The binding between curves is established by the Sigma protocol equality check, not by generator substitution. The Bulletproof transcript must include the BBS+ commitments to prevent splicing attacks.

\subsection{Epoch Nullifier Computation}

The client computes the epoch base point deterministically:
\[
H_T = \mathcal{H}_{\mathbb{G}_1}(\text{``EBUT:Epoch:''} \parallel T) \in \mathbb{G}_1.
\]
The hash function $\mathcal{H}_{\mathbb{G}_1}$ should be instantiated using a standard hash-to-curve method (e.g., RFC 9380) with domain separation tag \texttt{"ACT:Epoch:"}. The epoch $T$ is encoded as a 32-bit \textbf{little-endian} integer (matching the reference implementation's \texttt{t.to\_le\_bytes()} convention). The nullifier is $N_T = k_{\mathsf{sub}} \cdot H_T$. The server \textbf{must not} use a secret-derived base; the trustless nullifier property relies on the public derivation of $H_T$.

\subsection{Idempotency and Torn Transaction Recovery}
\label{sec:idempotency}

Network failures and client crashes can lead to retried requests. To prevent double-spend errors and improve user experience, the server must implement idempotency caches.

\paragraph{Refresh Phase Idempotency.}
When the server issues a Daily Token, it computes a cache key derived from the epoch nullifier $N_T$:
\[
\mathtt{key}_{\mathsf{ref}} = \mathcal{H}_{\mathsf{scalar}}\!\left(\text{``EBUT:Idempotency:Refresh''} \parallel N_T \right).
\]
The mapping $\mathtt{key}_{\mathsf{ref}} \mapsto (A_{\mathsf{daily}}, e_{\mathsf{daily}}, s'_{\mathsf{daily}})$ is stored for the remainder of the epoch (or at least several minutes). If a client resubmits a request carrying the same epoch nullifier $N_T$ (regardless of whether $K_{\mathsf{daily}}$ or $r_{\mathsf{daily}}$ was regenerated), the server returns the cached response \textbf{without} re-executing the expensive verification (pairings, MSMs) and \textbf{without} re-checking $\DB_{\mathsf{epoch}}$.

\textbf{Why the key must bind $N_T$, not $\pi_{\mathsf{refresh}}$.}
An earlier design keyed the cache on a hash of $\pi_{\mathsf{refresh}}$ directly.  This is insecure.  The proof $\pi_{\mathsf{refresh}}$ incorporates $K_{\mathsf{daily}}$, which in turn depends on the ephemeral blinding scalar $r_{\mathsf{daily}}$.  An adversarial client is not bound to reuse the same $r_{\mathsf{daily}}$ on a retry.  If the server crashes after writing the cache entry but before inserting $N_T$ into $\DB_{\mathsf{epoch}}$, an adversary can submit a new request with a freshly chosen $r_{\mathsf{daily}}' \neq r_{\mathsf{daily}}$, producing a different $K_{\mathsf{daily}}'$, a different $\pi_{\mathsf{refresh}}'$, and therefore a different cache key.  The new request misses the cache, passes the $\DB_{\mathsf{epoch}}$ freshness check (since the crash left $N_T$ uninserted), and the server issues a \emph{second} distinct daily token for the same epoch---shattering Rate-Limit Soundness.  Keying on $N_T$ closes this window: because $N_T = k_{\mathsf{sub}} \cdot H_T$ depends only on the client's master secret and the public epoch, an adversary cannot alter $N_T$ without forging a new master credential, and any retry with the same $N_T$ unconditionally hits the existing cache entry.

\textbf{Write ordering requirement.}
The idempotency cache entry \textbf{must} be written \emph{before} inserting $N_T$ into $\DB_{\mathsf{epoch}}$ (see Section~4.3.3, step~11). If the server crashes between the $\DB_{\mathsf{epoch}}$ insertion and the cache write, the client's retry will find $N_T$ already in $\DB_{\mathsf{epoch}}$ but no cache hit, and the server will permanently reject the user for the entire epoch. Writing the cache first eliminates this window: if the server crashes between the cache write and the $\DB_{\mathsf{epoch}}$ insertion, the retry finds a cache hit and returns the correct response without touching $\DB_{\mathsf{epoch}}$. If the idempotency cache and $\DB_{\mathsf{epoch}}$ reside in separate storage systems, these two writes must be wrapped in a single distributed transaction, or the cache must be written first.

\textbf{Crash recovery via deterministic nullifier derivation.}
Because $N_T = k_{\mathsf{sub}} \cdot H_T$ is deterministic, the $N_T$-keyed idempotency cache provides \emph{complete} crash recovery for honest clients.  If the client crashes after the server has processed the refresh but before the response is received, the client re-evaluates the PRF from its stored $k_{\mathsf{sub}}$, reconstructs the identical $N_T$, and the server's cache returns the cached blind signature.  No ephemeral state need be persisted beyond the long-lived master secret.  Importantly, even if the honest client re-derives fresh randomness $r_{\mathsf{daily}}'$ (producing a different $K_{\mathsf{daily}}'$ and $\pi_{\mathsf{refresh}}'$), the $N_T$-based cache lookup still hits the original entry and returns the cached daily token.

\paragraph{Spend Phase Idempotency.}
For spend requests, the cache key must bind the mode tag and every mode-specific public field using the same length-prefixed encoding:
\[
\begin{aligned}
h_{\mathsf{id}} = \mathcal{H}_{\mathsf{scalar}}\!\bigl(&\text{``EBUT:Idempotency:Spend''} \parallel \mathsf{mode} \\
&\parallel \mathsf{len}(s_{\mathsf{req}}^{\mathsf{enc}}) \parallel s_{\mathsf{req}}^{\mathsf{enc}} \\
&\parallel \mathsf{len}(C_s^{\mathsf{enc}}) \parallel C_s^{\mathsf{enc}} \\
&\parallel \mathsf{len}(C_{s-1}^{\mathsf{enc}}) \parallel C_{s-1}^{\mathsf{enc}} \\
&\parallel \mathsf{len}(\pi_{\text{BEQ\_s}}^{\mathsf{enc}}) \parallel \pi_{\text{BEQ\_s}}^{\mathsf{enc}} \\
&\parallel \mathsf{len}(\pi_{\text{BEQ\_s\_minus\_1}}^{\mathsf{enc}}) \parallel \pi_{\text{BEQ\_s\_minus\_1}}^{\mathsf{enc}} \\
&\parallel \mathsf{len}(\pi_{\mathsf{spend}}) \parallel \pi_{\mathsf{spend}} \parallel \eta \bigr).
\end{aligned}
\]
Here $s_{\mathsf{req}}^{\mathsf{enc}}$ is the 4-byte encoding of $s_{\mathsf{req}}$ in Exact mode and the empty string in Hidden mode; conversely $C_s^{\mathsf{enc}}$, $C_{s-1}^{\mathsf{enc}}$, $\pi_{\text{BEQ\_s}}^{\mathsf{enc}}$, and $\pi_{\text{BEQ\_s\_minus\_1}}^{\mathsf{enc}}$ are populated only in Hidden mode and encoded as the empty string in Exact mode. The fixed length of $\eta$ (16 bytes) eliminates ambiguity on that side; the length prefix on every optional field prevents byte-shifting and cross-mode malleability. Raw concatenation that omits the mode tag or omits the empty-field markers is \textbf{prohibited}: otherwise a relay can rewrite an Exact request as Hidden (or vice versa) while preserving the cache key. The server stores $h_{\mathsf{id}} \mapsto (A^*, e^*, s'^*)$ for the duration of the epoch plus the grace period---the same TTL assigned to the corresponding $k_{\mathsf{cur}}$ entry in $\DB_{\mathsf{spend}}$. \textbf{This TTL must exactly match or exceed the TTL of the token nullifier.} If the idempotency cache entry expires before $k_{\mathsf{cur}}$ expires in $\DB_{\mathsf{spend}}$, a client whose network dropped after the server committed the transaction but before the client received the response will later retry, find no cache hit, reach the double-spend check, and be permanently rejected with a 409 Conflict---losing their token with no recourse. \textbf{Crucially, the idempotency cache must be checked \emph{before} the double-spend database}, and must be \emph{written before} $\DB_{\mathsf{spend}}$ is mutated (see the exact spend flow and the Hidden-mode variant above). If a match is found, the cached refund token is returned immediately. This prevents legitimate retries from being incorrectly flagged as double spends. If the cache and the double-spend database reside in separate storage systems, these two writes must be wrapped in a single distributed transaction, or the cache must be written first.

\subsection{API Payload Specification}

Implementers must transmit the public statements alongside the proof. We provide JSON schemas for both refresh and spend requests. All binary data is encoded in hexadecimal.

\subsubsection{Refresh Request Payload}
\begin{verbatim}
{
  "T": <32-bit unsigned integer>,
  "N_T": <48-byte compressed G1 point, hex>,
  "K_daily": <48-byte compressed G1 point, hex>,
  "C_Delta": <48-byte compressed G1 point, hex>,
  "pi_refresh": {
    "A_prime": <48-byte compressed G1 point, hex>,
    "A_bar": <48-byte compressed G1 point, hex>,
    "T_BBS": <48-byte compressed G1 point, hex>,
    "T_scale_N": <48-byte compressed G1 point, hex>,
    "T_N": <48-byte compressed G1 point, hex>,
    "T_scale_D": <48-byte compressed G1 point, hex>,
    "T_D": <48-byte compressed G1 point, hex>,
    "T_scale_B": <48-byte compressed G1 point, hex>,
    "T_B": <48-byte compressed G1 point, hex>,
    "z_e": <32-byte scalar, hex>,
    "z_r1": <32-byte scalar, hex>,
    "z_s_tilde": <32-byte scalar, hex>,
    "z_k_tilde": <32-byte scalar, hex>,
    "z_c_tilde": <32-byte scalar, hex>,
    "z_E_tilde": <32-byte scalar, hex>,
    "z_u": <32-byte scalar, hex>,
    "z_v": <32-byte scalar, hex>,
    "z_w": <32-byte scalar, hex>,
    "pi_BEQ_exp": <binary blob, base64 or hex>
  }
}
\end{verbatim}

\subsubsection{Spend Request Payload}
The \texttt{eta} field carries the 128-bit session nonce generated locally by the client via CSPRNG (Section~\ref{sec:spend-exact}, step~1). It is encoded as 16 bytes in hexadecimal.

Exact mode:
\begin{verbatim}
{
  "mode": "Exact",
  "s": <32-bit unsigned integer>,
  "s_req": <32-bit unsigned integer>,
  "T_issue": <32-bit unsigned integer>,
  "k_cur": <32-byte scalar, hex>,
  "eta": <16-byte nonce, hex>,
  "K_prime": <48-byte compressed G1 point, hex>,
  "C_BP": <48-byte compressed G1 point, hex>,
  "pi_spend": {
    "A_prime": <48-byte compressed G1 point, hex>,
    "A_bar": <48-byte compressed G1 point, hex>,
    "T_BBS": <48-byte compressed G1 point, hex>,
    "T_scale_R": <48-byte compressed G1 point, hex>,
    "T_refund": <48-byte compressed G1 point, hex>,
    "T_scale_T": <48-byte compressed G1 point, hex>,
    "T_total": <48-byte compressed G1 point, hex>,
    "T_scale_BP": <48-byte compressed G1 point, hex>,
    "T_BP": <48-byte compressed G1 point, hex>,
    "z_e": <32-byte scalar, hex>,
    "z_r1": <32-byte scalar, hex>,
    "z_s_tilde": <32-byte scalar, hex>,
    "z_c_tilde": <32-byte scalar, hex>,
    "z_u": <32-byte scalar, hex>,
    "z_v": <32-byte scalar, hex>,
    "z_w": <32-byte scalar, hex>,
    "pi_BEQ_spend": <binary blob, base64 or hex>
  }
}
\end{verbatim}

Hidden mode:
\begin{verbatim}
{
  "mode": "Hidden",
  "T_issue": <32-bit unsigned integer>,
  "k_cur": <32-byte scalar, hex>,
  "eta": <16-byte nonce, hex>,
  "K_prime": <48-byte compressed G1 point, hex>,
  "C_BP": <48-byte compressed G1 point, hex>,
  "C_s": <48-byte compressed G1 point, hex>,
  "C_s_minus_1": <48-byte compressed G1 point, hex>,
  "pi_BEQ_s": <binary blob, base64 or hex>,
  "pi_BEQ_s_minus_1": <binary blob, base64 or hex>,
  "pi_spend": {
    "A_prime": <48-byte compressed G1 point, hex>,
    "A_bar": <48-byte compressed G1 point, hex>,
    "T_BBS": <48-byte compressed G1 point, hex>,
    "T_scale_S": <48-byte compressed G1 point, hex>,
    "T_S": <48-byte compressed G1 point, hex>,
    "T_scale_R": <48-byte compressed G1 point, hex>,
    "T_refund": <48-byte compressed G1 point, hex>,
    "T_scale_BP": <48-byte compressed G1 point, hex>,
    "T_BP": <48-byte compressed G1 point, hex>,
    "z_e": <32-byte scalar, hex>,
    "z_r1": <32-byte scalar, hex>,
    "z_s_tilde": <32-byte scalar, hex>,
    "z_c_tilde": <32-byte scalar, hex>,
    "z_amt_tilde": <32-byte scalar, hex>,
    "z_u": <32-byte scalar, hex>,
    "z_v": <32-byte scalar, hex>,
    "z_w": <32-byte scalar, hex>,
    "z_x": <32-byte scalar, hex>,
    "pi_BEQ_spend": <binary blob, base64 or hex>
  }
}
\end{verbatim}

\subsection{Serialization and Domain Separation}

\begin{itemize}[leftmargin=*]
\item \textbf{Group Elements:} All $\mathbb{G}_1$ elements must be serialized in \textbf{compressed} form (48 bytes for BLS12-381). Deserialization must validate that the point lies in the prime-order subgroup ($q \cdot P = 1_{\mathbb{G}_1}$) and is canonically encoded.
\item \textbf{Scalars:} All scalars must be serialized as \textbf{fixed-width 32-byte big-endian} integers. During deserialization, any 32-byte scalar $x$ must be rejected if $x \ge q$ (the group order). This prevents scalar malleability attacks.
\item \textbf{BEQ cross-curve integer $z$:} The BEQ response value $z$ is \textbf{not} a curve scalar; it is a plain unsigned integer in $[0, 2^{240}-1]$ and must be serialized as a \textbf{fixed-width 31-byte big-endian} unsigned integer.  \textbf{Implementations must deserialize $z$ as a raw integer and perform the bounds check $0 \le z \le 2^{240}-1$ before any other operation.}  Under no circumstances should $z$ be passed to a curve-scalar deserialization routine (e.g., \texttt{Scalar::from\_bytes\_mod\_order}) before this check: doing so silently applies modular reduction and will accept a maliciously crafted $z \ge q_{25519}$ that breaks the cross-curve binding (see the Verifier Bound Check in Section~\ref{sec:crosscurve} and Critical Component C2).
\item \textbf{Domain Separation:} All hash function invocations must include a domain separation label. In the reference implementation, the generator DST is \texttt{"ACT:BBS:"}, the epoch base-point DST is \texttt{"ACT:Epoch:"}, the BEQ challenge domain is \texttt{"ACT:BEQ:Challenge"}, and the Fiat--Shamir labels are \texttt{b"Refresh"} and \texttt{b"Spend"} respectively. Any interoperable implementation must use these exact strings. The context hash $H_{\mathsf{ctx}}$ (Section~4.1.1) must be included as the first element in every Fiat--Shamir challenge to bind all proofs to the deployment.
\item \textbf{Fiat--Shamir Hash Width:} For the BLS12-381 scalar field ($q$ a 255-bit prime), the bias introduced by reducing a 32-byte SHA-256 digest modulo $q$ is $(2^{256} \bmod q)/q \approx 2^{-255}$, which is negligible at the 128-bit security level. The reference implementation uses SHA-256 with repeated-subtraction mod-$q$ reduction and is cryptographically sound. For conservative interoperability or when targeting fields with smaller orders, use SHA-512 (64-byte output) before reduction.
\item \textbf{Subgroup Checks:} Every received $\mathbb{G}_1$ element must be validated to lie in the prime-order subgroup to prevent small-subgroup attacks.
\end{itemize}

\subsection{Nullifier Database Architecture}

The server maintains two logical sets:

\begin{itemize}
\item $\DB_{\mathsf{epoch}}$: stores compressed $\mathbb{G}_1$ points (48 bytes). Entries are assigned a TTL of \emph{epoch duration $+$ grace period} at insertion time. \textbf{The database must not be bulk-cleared at the epoch boundary}; individual entries expire after their TTL elapses.
\item $\DB_{\mathsf{spend}}$: stores canonical scalars (32 bytes). Entries are assigned a TTL of \emph{epoch duration $+$ grace period} at insertion time. \textbf{The database must not be bulk-cleared at the epoch boundary}; individual entries expire after their TTL elapses.
\item \textbf{Nonce store} ($\mathcal{N}_{\mathsf{spent}}$): stores the 16-byte session nonces $\eta$ that have been consumed during spend operations. Each entry \textbf{must} be assigned a TTL equal to \emph{epoch duration $+$ grace period}---the same TTL as $\DB_{\mathsf{spend}}$. Omitting a TTL causes permanent accumulation: at high spend velocity, uncontrolled growth of the nonce set will exhaust Redis memory. In Redis, store each nonce as \verb|SET "nonce:{eta_hex}" 1 EXAT {epoch_end + grace_seconds}| to guarantee per-entry expiry independent of other entries.
\end{itemize}

\paragraph{Epoch Grace Period.}
\textbf{Critical invariant:} the grace period and database expiry must be consistent. If the server accepts refresh or spend requests for epoch $T-1$ during a grace window of duration $\Delta_{\mathsf{grace}}$ after the start of epoch $T$, then entries written to $\DB_{\mathsf{epoch}}$ and $\DB_{\mathsf{spend}}$ during epoch $T-1$ \emph{must remain readable} for at least $\Delta_{\mathsf{grace}}$ into epoch $T$. Bulk-wiping either database at the epoch boundary while simultaneously accepting grace-period requests creates an exploit window in which an attacker may re-use an epoch-$T{-}1$ nullifier that appears fresh because its record was erased. Concretely, use per-key TTL-based storage: in Redis, store each nullifier as an individual key (for example, \verb|SET "epoch:nullifier:{N_T}" 1 EXAT {epoch_end + grace_seconds}|) so that each entry expires only after the grace window has fully closed, independent of other entries.

\paragraph{Atomicity.}
All insertions must be conditional and atomic. In SQL, use \texttt{INSERT} with \texttt{ON CONFLICT DO NOTHING}. In Redis, use \texttt{SETNX}. If an insertion fails due to a duplicate key, the request must be rejected immediately.

\paragraph{High-Velocity Spends.}
A user may legitimately generate up to $c_{\mathsf{max}}$ spends in a short period. To handle this efficiently, implementers \textbf{must} use an exact, authoritative data structure as the definitive store for $\DB_{\mathsf{spend}}$ and the nonce spent-set. In Redis, store each nullifier as an individual key with a per-key TTL (for example, \verb|SET "spend:nullifier:{k_cur}" 1 EX {ttl_seconds}|), which provides $O(1)$ membership tests and per-entry expiry without affecting other entries. A probabilistic structure such as a Bloom filter \textbf{must not} serve as the authoritative store: as the filter fills during an epoch, the false-positive rate rises, causing legitimate spends to be incorrectly rejected. A Bloom filter or Cuckoo filter \emph{may} be placed in front of the exact store as a fast-reject optimisation to avoid unnecessary database lookups for obviously duplicate entries, but any positive result from the filter must be confirmed against the authoritative store before the request is rejected.

\subsection{Verification Order (Optimized)}

To prevent CPU exhaustion attacks from malformed proofs, the server must perform checks in the following order (fastest to slowest). \textbf{All database writes (state mutations) are deferred to the final step and executed only after all cryptographic checks pass.} Performing writes earlier exposes the server to denial-of-service attacks in which an adversary submits a valid nonce with garbage cryptographic proofs, permanently burning the legitimate user's nonce before the proof is validated.

\begin{enumerate}[leftmargin=*]
\item Epoch and parameter bounds (e.g., $s > 0$, $T$ valid).
\item \textbf{Idempotency cache check} (return cached token if hit).
\item Nonce read-only check (Merkle proof if applicable; \textbf{no write}).
\item Nullifier read-only check and zero-scalar rejection ($N_T \neq 1_{\mathbb{G}_1}$, $k_{\mathsf{cur}} \neq 0$); \textbf{no write}.
\item Basic inequality checks ($A' \neq 1_{\mathbb{G}_1}$, $T_{\mathsf{BBS}} \neq 1_{\mathbb{G}_1}$).
\item Bridge equation checks (additive group operations, fast).
\item \textbf{Schnorr MSM verification} (BBS+ core without pairing) – $\approx 0.5$ ms.
\item \textbf{Bulletproof verification} (or BEQ verification) – $\approx 1$--$2$ ms.
\item \textbf{Bilinear pairing checks} – $\approx 2$--$4$ ms.
\item \textbf{Atomic state mutation} (write idempotency cache first; then insert nullifier into $\DB_{\mathsf{spend}}$, mark nonce as spent; reject on concurrent duplicate).
\end{enumerate}

\subsection{Edge Cases}

\begin{itemize}[leftmargin=*]
\item \textbf{$s = 0$:} Reject explicitly. Spending zero credits is meaningless.
\item \textbf{Zero scalars:} The client must ensure $k_{\mathsf{sub}} \neq 0$, $k_{\mathsf{daily}} \neq 0$, $k^* \neq 0$. The server must reject any nullifier $k_{\mathsf{cur}} = 0$ or $N_T = 1_{\mathbb{G}_1}$.
\item \textbf{Zero balance:} A token with $c_{\mathsf{bal}} = 0$ is allowed; further spends will fail because $s > 0$ and $m \ge 0$ cannot be satisfied.
\item \textbf{Clock synchronization:} Accept refresh requests with $T$ within a small tolerance (e.g., $\pm 1$ epoch) to accommodate client clock drift, but enforce the expiry check strictly.
\item \textbf{$c_{\mathsf{max}}$ range:} The server must enforce $c_{\mathsf{max}} \in [0, 2^{32}-1]$ during minting to ensure compatibility with 32-bit range proofs.
\end{itemize}

\subsection{Privacy and Operational Considerations}

\paragraph{Transport-Layer Anonymity via Local Nonce Generation.}
The session nonce $\eta$ is generated entirely on the client using a CSPRNG (Section~\ref{sec:spend-exact}, step~1), so no preliminary network request to the server is required.  This eliminates the IP-correlation vulnerability that arose from the former \texttt{GET /nonce} endpoint: a client that fetched a server-issued nonce and then submitted a spend request was trivially linkable at the network layer by any observer who could correlate the two requests.  By generating $\eta$ locally, the client makes only a single outbound connection per spend, natively preserving network-layer anonymity without any requirement for anonymous transports.

\textbf{Formal collision and replay resistance.}
The nonce $\eta \xleftarrow{\$} \{0,1\}^{128}$ is a fresh, locally sampled 128-bit random value.  The server maintains the exact spent-nonce set $\mathcal{N}_{\mathsf{spent}}$ and rejects any $\eta$ already present.  A replayed transcript is therefore rejected unconditionally.  For honest clients, the probability that two independently sampled nonces collide within the TTL grace window is bounded by the birthday paradox:
\[
p_{\mathsf{coll}} \;\le\; \frac{q_S^2}{2^{128}},
\]
where $q_S$ is the total number of spend requests within one epoch.  For any polynomial $q_S$, this probability is strictly negligible.

\paragraph{Epoch Boundary Burst Capacity.}
If $\Delta_{\mathsf{grace}} > 0$, a user may hold an unexpired token from the previous epoch while also minting a new token for the current epoch. During this overlap window, the user can spend up to $2 \cdot c_{\mathsf{max}}$ credits. This does not violate the long-term rate limit (one token per epoch), but it does permit a brief localized burst. If $\Delta_{\mathsf{grace}} = 0$ (the default), no such overlap exists. This trade-off is acceptable when fault tolerance matters more than strict instantaneous rate limiting; if strict per-minute rate limiting is required, the grace period must remain zero or be supplemented by an external system.

\paragraph{Cross-Epoch Linkability Assumption.}
The unlinkability of $N_T$ across epochs relies on the SXDH assumption in Type-3 pairings like BLS12-381. This provides computational (rather than information-theoretic) secrecy, which is the standard security guarantee for such protocols.

\paragraph{Token Theft Resistance.}
Even if an attacker intercepts a \texttt{SpendRequest} and submits it before the victim, the attacker receives only the blind signature $(A^*, e^*, s'^*)$. Without the local blinding factor $r^*$, the attacker cannot compute the final token secret $s^* = r^* + s'^*$. The token is cryptographically useless to the attacker, while the victim can safely resubmit (or receive from the idempotency cache) and unblind correctly.

\paragraph{Plaintext Spend Amount and Statistical Fingerprinting.}
In Exact mode (Section~\ref{sec:spend-exact}), the spend amount $s$ is transmitted in plaintext as a deliberately chosen application-layer design point: because the server already enforces a fixed price $s_{\mathsf{req}}$, the disclosure of $s = s_{\mathsf{req}}$ conveys no information beyond what the server already knows.  The resulting metadata channel—statistical fingerprinting of spend-amount sequences—is therefore classified as an \emph{out-of-scope application-layer policy assumption}, not a cryptographic break.  Deployments that require full amount privacy should instead use Hidden mode (Section~\ref{sec:spend-hidden}), which removes this metadata channel at the cost of larger proofs and additional verification work.

\section{Production Implementation: Critical, Medium-Priority, and Optimisation Decisions}
\label{sec:implementation-triage}
\label{sec:implementation}

The preceding sections present EBUT as a mathematical object.  A faithful,
production-grade instantiation must navigate a large design space in which
some decisions are \emph{cryptographically mandated} (omitting or weakening
them breaks the proven security properties), others are \emph{engineering
imperatives} (omitting them creates exploitable failure modes in real
deployments without affecting the formal proofs), and others are pure
\emph{optimisations} (their omission degrades performance but not security).
This section provides a unified, three-tier classification of every
significant decision in the reference implementation, and records the
precise ways in which the implementation diverges from the simplified
presentation in earlier sections.

\subsection{Classification Criteria}

\begin{itemize}[leftmargin=*]
\item \textbf{Critical~(C).} Omitting or weakening this component invalidates
  one or more of the formal security properties proved in
  Sections~\ref{sec:crosscurve} and~\ref{sec:security-ebut}: it either
  breaks soundness (unforgeability or rate-limit soundness), zero-knowledge
  (unlinkability), or completeness.
\item \textbf{Medium~(M).} The component does not change the cryptographic
  proof, but its omission creates a practical attack vector or permanently
  denies service to honest users under realistic network conditions.
\item \textbf{Optimisation~(O).} The component improves throughput, latency,
  or bandwidth without affecting correctness or security; it can be replaced
  by a simpler equivalent without breaking any proof.
\end{itemize}

\subsection{Critical Components}
\label{sec:critical}

\begin{enumerate}[leftmargin=*,label=\textbf{C\arabic*.}]

\item \textbf{Integer response $z_e$ computed over $\mathbb{Z}$, never reduced
  modulo a field order.}
  In the BEQ prover (Section~\ref{sec:crosscurve}), the response $z = u_m +
  c \cdot m$ is a plain 241-bit integer.  Reducing $z$ modulo $q_{\text{BLS}}$
  or $q_{25519}$ before transmission would silently truncate the high bits,
  causing the BLS and Ristretto verification equations to use inconsistent
  interpretations of $z$ and breaking either completeness or the extraction
  argument.  The implementation carries $z_e$ as raw little-endian bytes
  throughout; the two field-reduced values $z_m^{\text{BLS}}$ and
  $z_m^{25519}$ are derived independently inside the verifier by calling
  \texttt{scalar\_from\_le\_bytes\_mod\_order} and
  \texttt{DalekScalar::from\_bytes\_mod\_order} on the identical byte array.

\item \textbf{Bounds check $z_e < 2^{240}$ on the verifier side, performed on the raw integer before any scalar operation.}
  The verifier must reject any $z_e \ge 2^{240}$ \emph{and} must perform this check on the raw deserialized integer---not on a value that has already been passed through a curve-scalar routine.  An
  adversary who transmits $z_e \ge q_{25519}$ but less than $2^{241}$ will find that a Ristretto scalar library silently reduces $z_e$ modulo $q_{25519}$, making the Ristretto equation verify with a different
  effective integer than the BLS equation checks, breaking the
  cross-curve equality binding.  The check is strictly tighter than
  $z_e < 2^{256}$ because $q_{25519} \approx 2^{252}$ while $q_{\text{BLS}}
  \approx 2^{255}$.  Concretely: parse $z_e$ as a 31-byte big-endian unsigned integer, assert $z_e \le 2^{240}-1$, and only then derive $z_m^{\text{BLS}} = z_e \bmod q_{\text{BLS}}$ and $z_m^{25519} = z_e \bmod q_{25519}$ by calling \texttt{scalar\_from\_le\_bytes\_mod\_order} and \texttt{DalekScalar::from\_bytes\_mod\_order} on the raw bytes.

\item \textbf{240-bit blinding factor $u_m$ for BEQ statistical zero-knowledge.}
  The blinder $u_m$ must be drawn uniformly from $[0, 2^{240}-1]$.  The
  statistical distance between the real and simulated distributions of $z_e$
  is $c \cdot m / 2^{240} \le (2^{128}-1)(2^{32}-1)/2^{240} < 2^{-80}$
  (Theorem~5.3).  Using a shorter blinder (e.g., $u_m \in [0, 2^{128})$) would
  allow a computationally unbounded distinguisher to infer a lower bound on
  $m$ from $z_e$, breaking zero-knowledge.

\item \textbf{Subgroup membership check on every received $\mathbb{G}_1$ element.}
  Every compressed $\mathbb{G}_1$ point arriving over the wire must be
  deserialized via a routine that validates prime-order subgroup membership.
  A point in a small subgroup satisfies the pairing equation
  $e(A', \mathsf{pk}_S) = e(\bar{A}, g_2)$ with order dividing the cofactor,
  allowing a forgery that bypasses the BBS+ check entirely.  The
  implementation delegates this to \texttt{blstrs::G1Affine::from\_compressed},
  which performs cofactor-clearing and subgroup validation inside the
  assembly-optimised blst library.

\item \textbf{Zero-point and zero-scalar rejection.}
  The server must explicitly reject $A' = 1_{\mathbb{G}_1}$,
  $T_{\mathsf{BBS}} = 1_{\mathbb{G}_1}$, $N_T = 1_{\mathbb{G}_1}$, and
  $k_{\mathsf{cur}} = 0$.  A zero signature component ($A' = 1_{\mathbb{G}_1}$)
  satisfies the pairing equation trivially for any scalar responses.
  A zero epoch nullifier ($N_T = 1_{\mathbb{G}_1}$) provides no epoch binding.
  A zero token secret ($k_{\mathsf{cur}} = 0$) makes every derived token
  nullifier equal the same identity element, collapsing the double-spend
  protection.  The client must also ensure $k_{\mathsf{sub}} \neq 0$,
  $k_{\mathsf{daily}} \neq 0$, and $k^* \neq 0$ at key-generation time.

\item \textbf{TOCTOU ordering: nullifier written last, only after all proofs pass.}
  The nullifier must be inserted into $\DB_{\mathsf{epoch}}$ or
  $\DB_{\mathsf{spend}}$ only after every cryptographic check passes (see the
  refresh finalization step and the final state-mutation step of the spend
  flow).  Inserting the nullifier before verifying
  the ZK proofs creates a denial-of-service: an adversary submits a valid
  nullifier with a malformed proof, causing the server to burn the nullifier
  permanently while aborting during proof verification, locking out the
  legitimate owner for the remainder of the epoch.

\item \textbf{Write ordering: idempotency cache written before nullifier insertion; cache key must bind $N_T$.}
  If the server crashes after inserting $N_T$ into $\DB_{\mathsf{epoch}}$ but
  before writing the cached daily-token response, the client's retry finds no
  cache hit, hits the nullifier check, and is permanently locked out.
  Writing the idempotency entry \emph{first} eliminates this window: a crash
  before the nullifier insertion leaves the DB clean so the next retry
  succeeds via the cache.  This ordering is a protocol invariant, not an
  engineering preference: its violation causes an honest user to lose access
  without any cryptographic fault on their part.

  Furthermore, the idempotency cache key \textbf{must} be derived from the epoch nullifier $N_T$ (Section~\ref{sec:idempotency}) and not from $\pi_{\mathsf{refresh}}$ or $K_{\mathsf{daily}}$ alone.  The latter values incorporate the ephemeral blinding scalar $r_{\mathsf{daily}}$, which an adversarial client may re-sample between retries.  If the cache key were proof-based, an adversary who detects a crash (via timeout) could submit a second request with fresh randomness, producing a different cache key, missing the earlier entry, and---because $N_T$ was never inserted during the crash---obtaining a second distinct daily token for the same epoch.  An $N_T$-keyed cache closes this gap: the adversary cannot alter $N_T$ without forging a new master credential, so any request with the same $N_T$ necessarily hits the existing entry and receives the already-issued response.

\item \textbf{Deterministic derivation of $k_{\mathsf{daily}}$ and $r_{\mathsf{daily}}$ supports crash recovery; cache is keyed on $N_T$.}
  Because $k_{\mathsf{daily}}$ and $r_{\mathsf{daily}}$ are derived deterministically from the persisted master secret $k_{\mathsf{sub}}$ and the public epoch $T$ via the PRF (Section~4.3.1), a client that reboots after the server has processed the refresh but before the response arrived can re-evaluate the same PRF outputs from $k_{\mathsf{sub}}$ alone and reconstruct the identical $K_{\mathsf{daily}}$ and $\pi_{\mathsf{refresh}}$.  The server's idempotency cache, keyed on $N_T = k_{\mathsf{sub}} \cdot H_T$ (Section~\ref{sec:idempotency}), matches the resubmitted nullifier and returns the cached blind signature, which the client correctly unblinds.  No synchronous on-device write of ephemeral randomness is required before the network request: the only secret that must be durably stored is the long-lived $k_{\mathsf{sub}}$, which is already a baseline requirement for any resumable subscription.  Implementations that key the cache on $\pi_{\mathsf{refresh}}$ rather than $N_T$ are vulnerable to the double-mint attack described in Section~\ref{sec:idempotency} and must be updated.

\item \textbf{BEQ Bulletproof transcript binds BBS+ and bridge commitments.}
  The Merlin transcript for the Dalek range proof must be initialised with the
  BBS+ commitments ($A'$, $\bar{A}$, $T_{\mathsf{BBS}}$) and bridge
  commitments before the range proof is generated.  Without this binding the
  Bulletproof is transferable: an adversary can extract the range proof from
  one BBS+ session and replay it in a different session, breaking the
  unlinkability of the BEQ construction (Section~5.5).  The implementation
  appends each commitment via
  \texttt{transcript.append\_message(b"comm", \&comm\_bytes)} in the same
  order on both prover and verifier sides.

\item \textbf{Public, secret-free derivation of the epoch base point $H_T$.}
  The domain separation tag used to compute $H_T = \mathcal{H}_{\mathbb{G}_1}(
  \text{"ACT:Epoch:"} \parallel T)$ must not incorporate any server secret
  or session-specific value.  Any server-controlled secret in the DST makes
  $H_T$ server-specific, allowing a malicious server to equivocate across
  epochs and break the nullifier-binding argument in Game~$\mathsf{G}_4$ of
  the unforgeability proof.

\item \textbf{Variable-length fields must be length-prefixed in all hash inputs.}
  Every hash invocation that concatenates two or more variable-length byte
  strings must include a fixed-width length prefix for each variable-length
  field (Section~7.4).  Without length prefixes an adversary can shift bytes
  between adjacent fields to produce an identical hash input for two distinct
  proof statements.  Fixed-size fields (48-byte $\mathbb{G}_1$ elements,
  32-byte scalars, 4-byte integers) require no prefix.

\item \textbf{Domain separation at every hash call, with $H_{\mathsf{ctx}}$ first.}
  Every Fiat--Shamir challenge must include a protocol-specific domain label
  (e.g., \texttt{b"Refresh"} or \texttt{b"Spend"}) as its final input.
  Without domain separation a proof generated in the refresh phase could be
  re-interpreted as a spend proof and vice versa.  The context hash
  $H_{\mathsf{ctx}}$ must appear as the \emph{first} element of each challenge
  to prevent a prover from choosing parameters that produce a useful collision.

\item \textbf{Spend mode and mode-specific fields must be transcript-bound and syntactically enforced.}
  Exact and Hidden spend requests share a common proof skeleton.  If the
  verifier hashes only the common fields, a relay can rewrite the request mode,
  insert or remove $s_{\mathsf{req}}$, $C_s$, or $\pi_{\text{BEQ\_s}}$, and
  potentially obtain either a cross-mode replay or an idempotency-cache
  collision.  The Fiat--Shamir transcript and the spend idempotency key must
  therefore bind the mode tag plus a canonical encoding of every mode-specific
  field, including empty encodings for absent fields.  Exact mode must reject
  any non-empty Hidden-mode fields; Hidden mode must reject any disclosed $s$
  or $s_{\mathsf{req}}$.

\item \textbf{$H_{\mathsf{ctx}}$ binds all generators and both server keys.}
  The formula in Section~4.1.1 must bind not only the server public key but
  also both key types and all six $h_i$ generators (see updated formula in
  Section~4.1.1).  Including the generators prevents a parameter-substitution
  attack; including both $\mathsf{pk}_{\mathsf{master}}$ and
  $\mathsf{pk}_{\mathsf{daily}}$ prevents a token issued under one key being
  replayed against the other.

\item \textbf{Separate master and daily server keys.}
  The reference implementation uses distinct key pairs for master and daily
  token signing.  Although a single key pair suffices for the security proofs,
  key separation provides meaningful operational security: a compromise of
  the daily signing key does not invalidate long-term master subscriptions,
  and the daily key can be rotated at epoch boundaries.  Deployments should
  treat key separation as a production requirement.

\end{enumerate}

\subsection{Medium-Priority Components}
\label{sec:medium}

These components do not affect the formal proofs but are essential for a
secure, reliable production deployment.

\begin{enumerate}[leftmargin=*,label=\textbf{M\arabic*.}]

\item \textbf{Per-entry TTL expiry for nullifier databases; no bulk clear.}
  Entries in $\DB_{\mathsf{epoch}}$ and $\DB_{\mathsf{spend}}$ must expire
  individually via a per-key TTL of \emph{epoch duration $+$ grace period}.
  Bulk-wiping either database at the epoch boundary while simultaneously
  accepting grace-period requests opens a re-use window: a nullifier from
  epoch $T-1$ that was erased appears fresh, enabling a replay attack.

\item \textbf{Nonce store TTL matches $\DB_{\mathsf{spend}}$ TTL.}
  The 16-byte session nonce $\eta$ must be stored with the same TTL as the
  corresponding $k_{\mathsf{cur}}$ entry.  If the nonce store expires earlier,
  a client whose response was lost can resubmit after expiry, find no nonce
  hit, reach the double-spend check, and be permanently rejected.  Conversely,
  omitting the TTL entirely causes the nonce set to grow without bound.

\item \textbf{Idempotency cache TTL must not be shorter than the nullifier TTL.}
  As argued in Section~7.4, if the idempotency cache entry for a spend
  request expires before $k_{\mathsf{cur}}$ expires in $\DB_{\mathsf{spend}}$,
  a legitimate client retry is misclassified as a double-spend.  The TTLs of
  these two structures must be aligned.

\item \textbf{Verification order: cheapest checks first; all state mutations last.}
  Requests must be processed in the order: (1)~parameter bounds,
  (2)~idempotency cache, (3)~nonce/nullifier read-only checks, (4)~point
  validity, (5)~bridge equations, (6)~Schnorr MSM, (7)~BEQ range proof,
  (8)~bilinear pairing, (9)~atomic state mutation.  Performing a pairing
  check before cheaper structural checks wastes $\approx 3$~ms per
  malformed request; writing the nullifier before verification enables the
  TOCTOU attack described in~C6.

\item \textbf{Nullifier authoritative store must not use a probabilistic filter.}
  A Bloom filter or Cuckoo filter may serve as a fast-reject layer in front
  of the exact store, but the authoritative double-spend decision must use
  an exact data structure.  As the filter fills during an epoch its false-
  positive rate rises, causing legitimate spends to be rejected without
  cryptographic fault.

\item \textbf{Grace-period and DB TTL must be configured as a unit.}
  The invariant $\mathsf{TTL} \ge \mathsf{epoch\_duration} +
  \mathsf{grace\_period}$ must hold at all times.  Extending the grace window
  at runtime without also extending the DB TTL creates a re-use window.

\item \textbf{Nullifier insertion must be atomic and conditional.}
  The conditional insert must be a single atomic operation (e.g.,
  \texttt{INSERT~\ldots~ON~CONFLICT~DO~NOTHING} in SQL, \texttt{SETNX} in
  Redis).  A non-atomic read-then-write allows two concurrent requests with
  the same nullifier to both pass the read check before either write
  completes, enabling a double-spend.

\item \textbf{Transport-layer anonymity is natively preserved by local nonce generation.}
  Because the session nonce $\eta$ is generated locally on the client via a CSPRNG (Section~4.4.1), no preliminary network request to the server is required.  Each spend transaction requires only a single outbound connection, eliminating the IP-correlation window that existed when nonces were fetched from a server endpoint.  The EBUT cryptographic layer is therefore fully anonymous at the network layer by construction; no additional transport such as Tor is required for basic anonymity, though its use remains compatible and provides defence-in-depth.

\item \textbf{Epoch boundary burst capacity is an acceptable trade-off.}
  During the overlap of an active non-zero grace window a user may hold an unexpired
  token from epoch $T-1$ and a fresh token for epoch $T$, permitting a
  burst of up to $2 \cdot c_{\mathsf{max}}$ credits.  This does not break
  the long-term rate limit.  If strict per-epoch capacity is required the
  grace period must be set to zero.

\end{enumerate}

\subsection{Optimisation Components}
\label{sec:optimisations}

These components are pure performance improvements.  Each may be replaced by
a simpler, slower equivalent without affecting correctness or security.

\begin{enumerate}[leftmargin=*,label=\textbf{O\arabic*.}]

\item \textbf{ChaCha20Rng seeded from OS entropy for bulk blinder generation.}
  At the start of each \texttt{prove} call, 32 bytes are drawn from the
  caller's CSPRNG to seed a \texttt{ChaCha20Rng} instance; all blinder
  scalars are then generated from this fast SIMD stream cipher.  The 256-bit
  seed provides the full security level.  A compliant alternative is to call
  the OS RNG for each blinder directly, at a cost of $\approx 5$--$10\times$
  slower blinder generation.

\item \textbf{Fixed-base precomputed scalar multiplication (FixedBaseTable).}
  The seven protocol generators $g_1, h_0, \ldots, h_5$ are known at startup.
  For each, a $w=4$-bit windowed lookup table of $64 \times 15$ affine points
  ($\approx 90$~KB per table) is precomputed once.  Scalar multiplication
  then requires exactly 64 conditional mixed additions and zero doublings,
  compared with $\approx 256$ doublings and $\approx 128$ additions for
  double-and-add.  The simpler alternative is to use the library's standard
  variable-base scalar multiplication at no loss of security.

\item \textbf{$\mathbb{G}_2$Prepared pairing line-function caching.}
  The $\mathbb{G}_2$ elements $g_2$, $\mathsf{pk}_{\mathsf{master}}$, and
  $\mathsf{pk}_{\mathsf{daily}}$ never change after key generation.
  Pre-converting them to \texttt{G2Prepared} amortises the $\approx 3$~ms
  line-function computation across all verify calls.

\item \textbf{OnceLock singletons for Generators, BulletproofGens, and PedersenGens.}
  All three expensive-to-construct structures are initialised exactly once
  per process via Rust's \texttt{OnceLock} and shared across threads.  This
  eliminates $\approx 3$--$5$~ms of redundant hash-to-curve work per call.

\item \textbf{batch\_normalize: Montgomery batch field inversion.}
  Converting $N$ projective $\mathbb{G}_1$ points to affine via a naive loop
  requires $N$ field inversions.  \texttt{batch\_normalize} uses Montgomery's
  batch-inversion trick (one field inversion $+$ $3N$ multiplications),
  amortising the cost for any batch of points that must be normalised before
  the challenge hash.

\item \textbf{g1\_mul\_u32 for 32-bit scalar multiplication.}
  Spend amounts and epoch numbers are 32-bit values.  A 32-iteration
  double-and-add loop is used in place of a full 256-bit scalar
  multiplication, running $\approx 8\times$ faster with identical correctness.

\item \textbf{Schwartz--Zippel RLC: combined bridge and Schnorr MSM.}
  The multiple bridging-equality checks and the Schnorr equation are combined
  via a Schwartz--Zippel random linear combination with weights
  $c, c^2, c^3, 1$ derived from the Fiat--Shamir challenge, yielding a single
  Pippenger MSM.  Soundness error is $\le 4/|\mathbb{F}_r|$.  This reduces
  the number of MSMs from four to one, cutting the Schnorr/bridge portion of
  verification latency by $\approx 40\%$.

\item \textbf{Batch verification via RLC pairing aggregation and parallel BEQ.}
  \nolinkurl{verify_refresh_batch} and \nolinkurl{verify_spend_batch} aggregate
  $N$ proofs using per-proof RLC weights $\rho_i$: one Pippenger MSM for all
  Schnorr equations; one two-pair \nolinkurl{multi_miller_loop} for all
  pairing equations; and parallel \nolinkurl{rayon::par_iter} for the
  independent BEQ range proofs.  This reduces
  \nolinkurl{final_exponentiation} calls from $N$ to $1$.

\item \textbf{rayon::join for concurrent BEQ and pairing checks.}
  Within a single proof, the BEQ range-proof verification ($\approx 1$--$2$~ms)
  and the pairing check ($\approx 2$--$4$~ms) are mathematically independent.
  \texttt{rayon::join} executes them on separate threads, reducing combined
  latency from $\approx 5$~ms to $\approx 3$--$4$~ms.

\item \textbf{BEQ transcript compression: $R_{\text{BLS}}$ and $R_{25519}$ omitted.}
  The full proof tuple includes $R_{\text{BLS}}$ and $R_{25519}$.  The
  production wire format omits them (saving 80 bytes per proof) and stores
  only the 16-byte challenge $c$; the verifier reconstructs both $R$ values
  algebraically from the responses and re-hashes them to confirm the challenge
  (Section~5.2, Output item).

\end{enumerate}

\subsection{Deviations Between the Reference Implementation and the Protocol Specification}
\label{sec:deviations}

The reference implementation makes several choices that differ from or are
more conservative than the presentation in earlier sections.  Where they
conflict, the implementation is the authoritative reference.

\begin{itemize}[leftmargin=*]

\item \textbf{Domain separation strings.}
  Section~4.1.1 used \texttt{"EBUT:Generator:"} and \texttt{"EBUT:Context"}
  as example DSTs.  The actual implementation uses \texttt{"ACT:BBS:"} for
  generator derivation, \texttt{"ACT:Context"} for $H_{\mathsf{ctx}}$,
  \texttt{"ACT:Epoch:"} for the epoch base point, and
  \texttt{"ACT:BEQ:Challenge"} for the BEQ Fiat--Shamir challenge.  The
  Fiat--Shamir domain labels within the main proof are \texttt{b"Refresh"}
  and \texttt{b"Spend"}.  Any interoperable implementation must match these
  strings exactly.

\item \textbf{Epoch integer encoding: little-endian (not big-endian).}
  The epoch $T$ is encoded as a 32-bit \emph{little-endian} integer in the
  implementation (consistent with BLS12-381 scalar serialisation throughout
  the codebase).  All implementations must use little-endian encoding when
  computing $H_T = \mathcal{H}_{\mathbb{G}_1}(\text{"ACT:Epoch:"} \parallel
  T_{\text{LE}})$.

\item \textbf{SpendProof includes $T_{\mathsf{scale},R}$ and $T_{\mathsf{refund}}$.}
  Although the Exact-mode soundness argument can derive the refund bridge from
  the total bridge and the public relation $C_{\mathsf{total}} = K' + s h_4$,
  the canonical wire format retains $(T_{\mathsf{scale},R},
  T_{\mathsf{refund}})$.  The reference implementation includes both in the
  \texttt{SpendProof} struct and in the Fiat--Shamir challenge hash, which
  keeps the Exact and Hidden encodings closer and simplifies verifier code.
  Hidden mode requires the pair outright.

\item \textbf{Fiat--Shamir hash width.}
  Earlier drafts stated that \texttt{SHA-256~mod~q} is ``prohibited.''  For
  the BLS12-381 scalar field ($q$ a 255-bit prime), the bias from a 256-bit
  SHA-256 output reduced modulo $q$ is $(2^{256} \bmod q)/q \approx 2^{-255}$,
  negligible at the 128-bit security level.  The implementation uses SHA-256
  with repeated-subtraction mod-$q$ reduction and is sound.  The former
  prohibition is replaced by the clarification in Section~7.5.

\item \textbf{BEQ wire format.}
  The specification in Section~5.2 lists the canonical tuple
  \[
  \pi_{\text{BEQ}} = (C_{25519}, R_{\text{BLS}}, R_{25519}, z, z_{\text{BLS}}, z_{25519}, \pi_{\text{range}}).
  \]
  The production wire format is
  \[
  (c_{128}, C_{25519}, z_e, z_{r,\text{BLS}}, z_{r,25519}, \pi_{\text{range}}),
  \]
  omitting the $R$ points as described in optimisation~O10.

\end{itemize}

\section{Performance Evaluation}
\label{sec:performance}

We report measured performance of the reference Rust implementation on BLS12-381
and Curve25519 (Ristretto).  All timings are averages of 10 iterations in
release mode (with \nolinkurl{RUSTFLAGS="-C target-cpu=native"}, enabling AVX2
acceleration in the Dalek Bulletproof library), on a single thread unless stated
otherwise.  Min/max across the 10 iterations are reported in the component
breakdown (Table~\ref{tab:components}) to characterise JIT warm-up effects.
The benchmark source is \nolinkurl{tests/performance.rs}.

\subsection{Proof Sizes}

\begin{table}[htbp]
\centering
\small
\begin{tabular}{|l|r|l|}
\hline
\textbf{Phase} & \textbf{Wire size} & \textbf{Breakdown} \\
\hline
Master Mint request  & 160 B & $K_{\mathsf{sub}}$ (48 B G1) + $T_{\mathsf{PoK}}$ (48 B G1) + $s_k, s_r$ (2$\times$32 B Fr) \\
Master Mint response & 112 B & $A_{\mathsf{sub}}$ (48 B G1) + $e, s'$ (2$\times$32 B Fr) \\
\hline
Epoch Refresh proof  & \textbf{1620 B} & BBS+ core $A', \bar{A}, T_{\mathsf{BBS}}$: 144 B (3$\times$G1) \\
                     &       & Bridge commitments: 288 B (6$\times$G1) \\
                     &       & Public values $N_T, K_{\mathsf{daily}}, C_\Delta$: 144 B (3$\times$G1) \\
                     &       & Schnorr responses: 288 B (9$\times$Fr) \\
                     &       & BEQ proof (dual-curve $+$ Bulletproof): 756 B \\
\hline
Spend proof (Exact mode) & \textbf{1548 B} & BBS+ core: 144 B (3$\times$G1) \\
                     &       & Bridge commitments: 288 B (6$\times$G1) \\
                     &       & Public values $K', C_{\mathsf{BP}}$: 96 B (2$\times$G1) \\
                     &       & Disclosed ($s$, $T_{\mathsf{issue}}$): 8 B (2$\times$u32) \\
                     &       & Nullifier $k_{\mathsf{cur}}$: 32 B (Fr) \\
                     &       & Schnorr responses: 224 B (7$\times$Fr) \\
                     &       & BEQ proof: 756 B \\
\hline
\end{tabular}
\caption{Wire-format proof sizes for each EBUT phase.  Spend figures are for Exact mode.}
\label{tab:sizes}
\end{table}

The BEQ proof of 756 bytes breaks down as: 16 B challenge $c_{128}$,
32 B Ristretto commitment $C_{25519}$, 32 B integer response $z_e$,
32 B BLS12-381 randomness response $z_{r,\mathsf{BLS}}$, 32 B Ristretto
randomness response $z_{r,25519}$, 4 B length prefix, and 608 B for the
32-bit Dalek Bulletproof.  The Schnorr commitment pair $(R_{\mathsf{BLS}},
R_{25519})$ is omitted from the wire format via transcript compression
(Section~8, O10) saving 80 bytes per proof.

All reported Spend proof sizes in Table~\ref{tab:sizes} are for Exact mode.
Hidden mode adds the public amount commitment $C_s$, a second BEQ proof
$\pi_{\text{BEQ\_s}}$, and hidden-mode-only bridge commitments and responses,
so its wire size is necessarily larger.

\subsection{End-to-End Timings}

\begin{table}[htbp]
\centering
\small
\begin{tabular}{|l|r|r|r|}
\hline
\textbf{Operation} & \textbf{Avg (ms)} & \textbf{Wire (B)} & \textbf{Ops/s} \\
\hline
Setup: generator creation (one-time) & 96.5 & --- & 10.4 \\
Server key generation                &  1.8 & --- & --- \\
$H_{\mathsf{ctx}}$ computation       &  0.5 & --- & --- \\
\hline
Master Mint: client begin (commit + PoK prove)    & 1.5 & 160 & 651 \\
Master Mint: server issue (PoK verify + BBS+ sign)& 2.5 & 112 & 405 \\
Master Mint: client finalize (unblind)            & $<$0.001 & --- & --- \\
\hline
Epoch Refresh: client prove  & 19.1 & 1620 &  52 \\
Epoch Refresh: server verify &  8.1 & 1620 & 124 \\
\hline
Spend (Exact): client prove  & 16.2 & 1548 &  62 \\
Spend (Exact): server verify &  6.4 & 1548 & 156 \\
\hline
\end{tabular}
\caption{End-to-end timings and single-threaded throughput.  Spend figures are
for Exact mode.  Measured at 10 iterations each, release build with AVX2 enabled.}
\label{tab:e2e}
\end{table}

Generator creation is a one-time startup cost.  All subsequent protocol
operations use the cached \texttt{Generators} singleton; the 96.5~ms cost is
not incurred on a per-request basis.

\subsection{Component Breakdown}

\begin{table}[htbp]
\centering
\small
\begin{tabular}{|l|r|r|r|}
\hline
\textbf{Component} & \textbf{Avg (ms)} & \textbf{Min (ms)} & \textbf{Max (ms)} \\
\hline
Single pairing $e(G_1, G_2)$ [baseline]     & 1.47 & 1.23 & 1.82 \\
Multi-pairing 2-input [used in verify]       & 2.20 & 1.68 & 2.65 \\
\quad Miller loop only (1 pair)              & 0.56 & 0.37 & 0.70 \\
\quad Final exponentiation                   & 1.09 & 0.75 & 1.25 \\
\hline
MSM  3 bases (issuance)                      & 0.72 & 0.49 & 0.97 \\
MSM  5 bases (BBS+ Schnorr check)            & 1.23 & 0.79 & 1.41 \\
MSM 12 bases (bridge + BBS combined)         & 1.69 & 1.17 & 2.53 \\
\hline
Hash-to-G1 ($H_T$ epoch computation)         & 0.147 & 0.125 & 0.161 \\
BEQ prove (dual-curve, 32-bit range)         & 9.67 & 8.75 & 11.55 \\
BEQ verify (dual-curve, 32-bit range)        & 2.43 & 1.98 & 3.35 \\
G1 scalar-mul (BBS+ randomisation $A\cdot r_1$) & 0.202 & 0.159 & 0.251 \\
SHA-256 challenge hash ($\approx$1~KiB input)& 0.012 & 0.008 & 0.040 \\
\hline
\end{tabular}
\caption{Per-component timings (10 iterations, release + AVX2).  The
multi-pairing check combines two Miller loops before a single
final exponentiation, saving one $\approx$1.1~ms final-exp versus two
independent pairings.}
\label{tab:components}
\end{table}

\subsection{Throughput and Latency Analysis}

\paragraph{Server-side throughput (single-threaded).}
At 124~ops/s for Epoch Refresh verification and 156~ops/s for Exact-mode Spend
verification, a single CPU core can handle 69 combined refresh+spend
operations per second.  Because BEQ verification and the pairing check
are executed concurrently via \texttt{rayon::join} (Section~8, O9), the
effective server latency per proof is approximately 6--8~ms rather than
the 9--10~ms that sequential execution would require.  Multi-core scaling
is nearly linear for independent proofs via the batch verification paths
\texttt{verify\_refresh\_batch} and \texttt{verify\_spend\_batch}
(Section~8, O8), which aggregate $N$ proofs into one Pippenger MSM and
one final exponentiation regardless of $N$.

\paragraph{Client-side latency.}
Epoch Refresh proof generation takes $\approx$19~ms and Exact-mode Spend proof
generation takes $\approx$16~ms on the same hardware.  The dominant cost
in both is the Dalek Bulletproof generation inside the BEQ (9.7~ms avg),
which runs on Curve25519 and benefits from AVX2 vectorisation.  The
BBS+ randomisation and Schnorr blinder generation together cost
$\approx$2~ms; the Fiat--Shamir challenge hash is negligible at 0.012~ms.
Hidden mode is strictly slower because it adds a second BEQ proof and extra
bridge terms for the hidden amount commitment.

\paragraph{Proof sizes.}
Each full Refresh proof and Exact-mode Spend proof fits in well under 2~KB
(1620~B and 1548~B respectively), making the protocol suitable for TLS-protected
HTTP/2 connections without fragmentation.  The BEQ component accounts
for 756~B of each proof; the transcript compression described in
Section~8 (O10) already saves 80~B per proof compared to the
uncompressed format.  Hidden mode is larger because it carries an additional
amount commitment and BEQ proof.

\paragraph{Comparison with earlier estimate.}
The earlier draft estimated Refresh proof size at $\approx$3.0~KB and
Exact-mode Spend at $\approx$3.2~KB.  The reference implementation achieves
1.62~KB and 1.55~KB respectively—a reduction of approximately 47\%—
due to the dual-curve BEQ design (which produces a compact 756~B
combined Sigma + Bulletproof instead of a standalone BLS12-381 range
proof), and the transcript compression that omits the Schnorr
commitment pair from the wire format.  Hidden mode intentionally trades some
of that compactness for amount privacy.

\subsection{Applied Optimisations}

The following optimisations are present in the reference implementation
and reflected in the measurements above; full classification and
rationale appear in Section~\ref{sec:optimisations}.

\begin{enumerate}[leftmargin=*]
\item \textbf{Pippenger MSM} for all multi-scalar multiplications
  ($\geq 12$ bases per verify call).  Compared to a naive double-and-add
  fold the measured speedup for the 12-base bridge+BBS combined check is
  $\approx$5--10$\times$.

\item \textbf{Single multi-pairing} replaces two independent pairing
  checks.  Two Miller loops are combined before one final exponentiation,
  reducing the measured pairing cost from $2\times 1.47$~ms = 2.94~ms
  to 2.20~ms ($\approx$25\% saving).

\item \textbf{$\mathbb{G}_2$Prepared line-function caching} for $g_2$,
  $\mathsf{pk}_{\mathsf{master}}$, and $\mathsf{pk}_{\mathsf{daily}}$.
  Pre-conversion to \texttt{G2Prepared} amortises $\approx$3~ms of
  line-function computation across all verify calls.

\item \textbf{Dual-curve BEQ} replaces BLS12-381-native Bulletproofs for
  all range proofs.  The Dalek library's AVX2-optimised inner-product
  argument on Curve25519 is $\approx$3$\times$ faster than the equivalent
  BLS12-381 range proof, and the 32-bit Dalek Bulletproof (608~B) is
  smaller than a BLS12-381 256-bit range proof.

\item \textbf{\texttt{rayon::join}} for concurrent BEQ + pairing
  verification within a single proof call, reducing combined latency
  from $\approx$10~ms to $\approx$6--8~ms.

\item \textbf{OnceLock singletons} for \texttt{Generators},
  \texttt{BulletproofGens}, and \texttt{PedersenGens} eliminate
  3--5~ms of redundant hash-to-curve work per call.
\end{enumerate}

\section{Conclusion}

We have presented Epoch-Bound Usage Tokens (EBUT), a primitive that enforces rate-limited anonymous access by cryptographically binding tokens to discrete time epochs. Our construction combines blind BBS+ signatures, Pedersen commitments, and zero-knowledge proofs to achieve strong privacy guarantees while preventing token accumulation and burst consumption. We formalized the security properties of EBUT—unforgeability, unlinkability, and rate-limit soundness—and provided rigorous, game-based proofs under standard assumptions. The integration of cross-curve batched equality proofs ensures practical efficiency without compromising security. We provided a complete, production-ready specification with detailed implementation guidance, API schemas, and performance estimates. EBUT offers a robust solution for privacy-preserving subscription services and rate-limited resource access in anonymous systems.

\bibliographystyle{alpha}
\begin{thebibliography}{99}

\bibitem{act}
Anonymous Credit Tokens (ACT): A Complete Formal Specification.
\emph{Cryptology ePrint Archive}, 2024.

\bibitem{NTAT}
F.~B.~Durak, A.~Talayhan, L.~Marco, and S.~Vaudenay.
Non-Transferable Anonymous Tokens by Secret Binding.
\emph{Full version, May 8th, 2024}.

\bibitem{CamenischL01}
J.~Camenisch and A.~Lysyanskaya.
An Efficient System for Non-transferable Anonymous Credentials with Optional Anonymity Revocation.
\emph{EUROCRYPT 2001}.

\bibitem{CL02}
J.~Camenisch and A.~Lysyanskaya.
A Signature Scheme with Efficient Protocols.
\emph{SCN 2002}.

\bibitem{BBS04}
M.~H.~Au, W.~Susilo, and Y.~Mu.
Constant-Size Dynamic $k$-TAA.
\emph{SCN 2006}.

\bibitem{CFRG-BBS}
M.~Lodder, T.~Looker, and A.~Whitehead.
The BBS Signature Scheme.
\emph{IETF Internet-Draft draft-irtf-cfrg-bbs-signatures-05}, 2024.

\bibitem{BB04}
D.~Boneh and X.~Boyen.
Short Signatures Without Random Oracles.
\emph{EUROCRYPT 2004}.

\bibitem{Bulletproofs}
B.~B{\"u}nz, J.~Bootle, D.~Boneh, A.~Poelstra, P.~Wuille, and G.~Maxwell.
Bulletproofs: Short Proofs for Confidential Transactions and More.
\emph{IEEE S\&P 2018}.

\bibitem{PointchevalStern00}
D.~Pointcheval and J.~Stern.
Security Arguments for Digital Signatures and Blind Signatures.
\emph{Journal of Cryptology}, 13(3):361--396, 2000.

\bibitem{BonehShoup}
D.~Boneh and V.~Shoup.
A Graduate Course in Applied Cryptography.
\emph{Version 0.6}, 2023.

\end{thebibliography}

\end{document}
